\chapter{气体动力学基础}
本章应用理想流体运动基本方程研究可压缩流体的运动规律，也就是研究流体 密度有较大变化、热力学过程和力学过程捐合的流动规律.

\section*{7.1  气体动力学基本方程}\addcontentsline{toc}{section}{7.1 气体动力学基本方程}

当流体质点在运动过程中，密度发生较大变化时，需要考虑流体压缩性的影响. 气体是可压缩流体，气体动力学属于可压缩流体动力学. 气体动力学研究范围很广，包括高速空气动力学、气体波动力学和高温气体力学 等.本章作为气体动力学的基础，主要研究完全气体运动中热力学过程与力学过程捐 合时出现的现象，对它们进行分析，并给出简单的计算方法. 1.理想完全气体模型 本章作为气体动力学的基础知识，假定流体是理想的完全气体，即动力粘性系数 $\mu$=0. 热传导系数 ${\rm A}$=O. 应力张量飞= $\rho$lJ'j' 理想完全气体满足以下状态方程 z

\begin{equation*}\rho  = RpT\tag{7.1a}\end{equation*}

它是常比热容气体，其内能 r 、蜡 h 和情 5 分别为

\begin{gather*}
h = c\theta T\\
e = cvT
\end{gather*}

(7.1b) (7.1c)

s = c\textbackslash\{\}， ln 乓p.

\[R = c\theta  - Cv\]

(7.1d) (7.1e) 8341 R= 一一一 (7. lf) 式中 R 是气体常数， $\mu$ 是相对分子质量，如空气的相对分子质量是 29 ， R=287N. m/ (kg. K) ，勺 ， Cv 分别为比定压热容和比定容热容 ， Y=C${\rm þ}$/Cv 称为热容比或绝热指数. 在常温常压下的一般气体，如空气、氧气、氢气等，都可近似为理想完全气体. 2. 气体动力学基本方程 应用第 4 章的流体力学基本方程，气体运动应满足以下方程 z (1)质量守恒方程或连续方程 EE+v-dl==0(7.2a) rJt ${\rm •}$ 或将式 (7.2a) 中散度项展开 V ${\rm •}$ rJ.l=pV ${\rm •}$ U+U ${\rm •}$ 峙，代人式 (7.2a) 后得以下形式 的连续方程 z

\[Dp+oV.U=O (7.2b) Dt .\]

(2) 动量方程或运动方程 z 就是第 5 章导出的欧拉方程，但是本章讨论的流动过 程中密度 p 是变量 z 学 +U-VEI=-l vp(7.3)

\[rll\rho \]

(3) 能量方程 z 将理想流体的应力状态代人第 4 章的能量方程，可得 I ${\rm •}$ U 2 、 I ， U 2 、 1\textasciitilde{}， I e + \textasciitilde{}'> 1+ U ${\rm •}$ V I e + \textasciitilde{}'> 1 =一$\tilde{V}$. (pU) +q 飞 2 ，飞 2 , p 式中 U 2 =IU 尸，将方程右端的散度展开 V ${\rm •}$ ($\rho$U)=pV ${\rm •}$ U+U. V$\rho$ ，并将内能 e=h -$\rho$ /p 代人上式，经简单代数运算后，可得以下形式的能量方程 z 1 , , U2 \textbackslash\{\} , rT \textasciitilde{} 1 , , U2 \textbackslash\{\} 1 rJ /J :: (h + $\tau$1 + U ${\rm •}$ V ( h + $\tilde{n}$ 1 = \textasciitilde{} ':Y + ci (7. 4) 飞 2 ，飞 2' P rJ t 骂 式中 q = Dq/Dt 是单位质量气体质点吸人的热量.若流动是绝热 (q =0) 且连续的，即 过程是绝热可逆的，则由热力学第二定理 : Ds=Dq/T ， 可导出蜻增 i=EE= 旦 =0 哩 Dt T 就是说绝热连续的流动过程是等蛐过程 ， s= 常数，对于完全气体 s=cvln(p/$\rho$$\eta$ ，因 此有 旦/主\textbackslash\{\}= 0 Dl 飞 p1 , - 就是说，绝热的理想气体连续运动中，质点的情沿轨迹不变. (7.5)

\section*{7.2  声传播方程和马赫数}\addcontentsline{toc}{section}{7.2 声传播方程和马赫数}

一般情况下由式 (7.1a) ，式 (7.2a) ，式 (7. 3) ，式 (7 . 4) 共 6 个方程可解出流场物 理量 p ， 户 ， U ， T. 完全气体的温度则由状态分程求出.对于等情流动(连续绝热流动) , 流场物理量 p' $\rho$ ， U 可由方程组 (7.2a) ，式 (7. 3) ，式 (7. 5) 解出.从理想完全气体运动 的基本方程来看，可压缩流动的主要特点是有密度和其他热力学状态参数，如温度、 热始以及摘的变化.

1.气体的一维非定常微扰流动和声速 从物理学中我们知道，声音以波的形式传播，它是一种微弱的气体披运动，通常 声波在 三 维空间传播，为了便于理解，先讨论平面声波的流场. 图 7. 1 所示 一 个半元限长的绝热等截面管道，左端有一振动膜(相当于一个扬声 器) ，右端为无限长.初始状态等截面管道内的气体处于静止平衡状态 (U=O) ， 气体 的状态参数为如，向 $\cdot $ 在薄膜激振下，管内气体开始运动.讨论气体的微扰运动，也就 是说振动膜的振幅和运动速度很小，因此，它所激起的气体运动速度 u' (.1", t) 很小，运 动气体的状态参数 $\rho$ ， p 偏离原来的平衡值扣， $\rho$。 也很小. 微幅振动 薄膜 /',/ "气//.//.//,<, '//'//////// 图 7 . 1 推导声速用图 在等截面管道中，可以假定 气 体流动是 一 维的，即气体的速度、压强和密度等都 是 (.r， t) 的函数，即 u= u( .r， 仆， $\rho$=$\rho$ ( .r .t) ， p=p( .r， t) ， 由气体运动的基本方程，可以 得到以下简化的气体 一 维非定常运动的基本方程. 连续方程 运动方程 经+但主 =0 ${\rm e}$J t J .1"

\[Ju , dU 1 eJ/J\]

一十 H 一一一

\[eJ l ' - a.r p ().r\]

(7.6) (7. 7) 由于声波是绝热微扰运动，它可近似为等蛐过程，即有方程 z Ds =旦\{主 )+u 立\{主 )=0 (7.8) Dl Jt \textbackslash\{\}pY J $\theta$z 飞 pY J 假定初始时刻管内气体是静止且热力学平衡的，即 t=O 时: u( x ,O) =0 ， $\rho$ (x ， O) =

Po ,p (.r, 0) =po ,s( .r, O) =So. 由于情沿轨迹不变，故由式 (7.8) 可导巾，流场始终保持 均惰 ， s( .r， t)=so' 或写成 主=主o p1 pr, 当流场中受到微小激振后，扰动将在流场中传播，受扰动的流场可表示为

\[\rho =\rho .  +\rho '\cdot p=\rho  +p' ,  u=u'\]

(7.9) 已假定扰动是微弱的，即相对值 $\rho$'/$\rho$ ，， -O($\epsilon$ )<< l. p' /向 -O(d<<l ，户 ， U'2/ Po -O(d << 1.将其代入式 (7.6)- 式 (7.9) ，并将导数展开.就有 罕+向毛 +24旦;二 o

\[rll (/J. . rl I'\]

旦旦:牛 H' 豆豆=一 -L E正如\_ \_!\_\_(1 - e\_ \textbackslash\{\}'!\_主\textasciitilde{} = \_!\_\_ ( 1 - p' , (也还 Jt ,- J.l' 户 + p' J .r -如 $\gamma$pn I J.r po 飞- po 1\textbackslash\{\} dp 人 J .r 略掉二阶以上的小量可使方程线性化为 Jp' 1 ${\rm e}$1 u' \_ " Jt ' ${\rm o}$ V rJ.r 年:+豆豆 =0 dt 向 rJ:I ' (7.10) (7.11) 式中 c\textasciitilde{} = (d$\rho$ /dp) ， = r$\rho$。/向是等蛐过程中压强对密度的导数.由式( 7. 10) 和 式 (7.11) 分别消去 p' 或 u' 可得 旦旦:-JE正 -(1 Jti

\begin{gather*}
< (1\\
rJ.r2 - v
\end{gather*}

当:-t i 当 =0 J( a .r 二 (7.12) (7.13) 式 (7.12) 和式 (7.13) 是气体微扰运动所满足的方程，也称作一维声波方程.这是典型 的双曲型方程，其解的一般形式为 u=fl (.r -c"t)+I (.r十 Cll t) (7. 14) 函数扎 .I - 由初始条件确定.由解式 (7.14) 可以看出微小扰动在静止完全气体巾传 播有以下特点(参见图 7.2 )， I1 10 X

\[./; (x-c ,,t l\]

图7.2

\[L(x+c ,,t l\]

声传播的特征

(1)方程的解是两族简单波的叠加 p (2) 右传波 f$\rightarrow$ (:r - c" t) : 函数 J ， 沿迹线 :r - c" t = const. 不变.传播速度 d :r/

\[dt = cu ;\]

(3 )左传波 J- (:r + c" t ) :雨数 f 沿迹线.r + Co t = const. 不变，传播速度 d.r/

\[dt = -c,, ;\]

(4) 扰动传播速度 c o= j${\rm e}$布及万-:: = Ji${\rm i}$;，石7 . 这就是声波的传播速度.简称声 速.声速只和气体的热力学状态有关，与扰动的运动学特性，如扰动的频率、波长等 无关. 因此等截面绝热管道中 气 体的 一 维声传播是非色散性的政向波.以上分析可以推 广到声波在无界空间巾的传播，其结论是任意方向的 声 传播都是非色散的等速行进波. 2. 声速、马赫数 (1)完全气体的声速 上面讨论了气体平面运动中微弱扰动以波传播的形式运动，在一般均匀静止空 间的气体中，声波以球面波的形式传播，用同样的方法可导出等崎传播的声速公式: c2 = (i主) (7.15) 飞 J$\rho$/ 、 对于完全气体，将等摘关系式 $\rho$/l = const. 代人，可得声速公式: c = JYP = .j页T (7.16) v p 所以完全 气 体声速 c 只是温度的函数.在非均匀的流场中，不同时刻、不同点上 声速大小和当时当地的温度有关，温度越高，声速越大. (2) 马赫数 定义 7.1 流体速度 u 与当地声速 r 之比称为马赫数(用 Ma 表示): Ma=u /c. 马赫数的物理意义可解释为: ${\rm 1}$ Ma 是单位质量流体的惯性力与压强合力的量级之比，可做以下的量级估计: 惯性力 =止豆豆丛L上 \textasciitilde{}旦旦 L \_ vMn 2 压强合力 I V$\rho$ / p . I $\rho$ / Idp y … ${\rm 2}$ Ma 是气体质点单位质量的功能与内能的量级之比 z 动能 \_ U2 /2 U2/2 \_ y(y 一 1)施 8 . . 内能 L\textbackslash\{\}.. T $\rho$ / (y 一 1) p 2" 白 另外，我们知道 ( 2 = (dp / dp) ${\rm ‘}$表示流体的压缩性. (dp / dp) ， 越大，表示压缩性越 小，而声速大 ;(dp/dp) ，越小，反映流体的压缩性大，但声速小.例如，理想的不可压 缩流体，它的声速是无限大，任何小扰动在 一 瞬间就传播到全流场.从马赫数的定义， 可以看到如果流体运动速度相同.则流体的声速越大，马赫数越小，它反映这种流体 运动的压缩性也就越弱.例如:理想不可压缩流体的马赫数趋向于零 z 常温 (293K)

下空气的声速约为 340m / 日，以每小时 120km( 33. 3m / s) 行驶的汽车，它的运动马赫 数约为 0.1; 以每小时 1000km(277.8m / s) 飞行的喷气飞机，它的飞行马赫数约为 0.82. 由估算的马赫数，可以判断，喷气飞机周围流动的压缩性大于汽车周围的气体 流动的压缩性. ${\rm 3}$气体流动分类 从马赫数的定义可知 : Ma=l 表明流体质点速度与当地声速相同，称这种流动 状态为临界状态.马赫数大于 1 或小于 l 的流动分别称作起声速流动或亚声速流动. 在空气动力学中通常把气体流动分成以下四类: 岛公 < 1 亚声速流动 Ma 句 1 跨声速流动 Ma > 1 超声速流动 J\textbackslash\{\}.也 >> 1 高超声遥流动 3. 超声速 (Ma>l) 流场与亚声速 (Ma < l) 流场的主要差别:影响 域和依赖域 静止气体中的运动物体或均匀气流绕物体产生的流场都可以看作是气体连续受 扰运动.超声速或亚声速运动的物体都能产生气体扰动，但是扰动在气体中的传播有根本 性差别.为了考察超声速流场和亚声速流场的最主要的差别，假设微小扰动只是一个点. 定义 7.2 点扰动能够传播到的空间区域.称为扰动的影响域. 下面来考察亚声速或超声速运动的点扰动的影响域. (1)静止流场中固定的点扰动源(坐标原点)的影响域.此时扰动在各方向以相 同的声速 c 传播， 1对此静止点扰动以同心阅球的披面向四周传播，见图 7.3(a). 例如， 从 t=O 发出的扰动，在 t = 1 时.该扰动传播到半径为 ct 的球面。 l 上 ;t = 2 时，该扰 动传播到半径为 2ct 的球面 O 2 上 ;t=3 时，该扰动传播到 3ct 的球面0;.上;依次类 推，足够长时间后，扰动将传播到全流场.也就是说，影响域是全部空间. (2) 亚声速运动的点扰动的影响域(图 7.3(b)). 此时扰动点运动速度 U 小于声 速 c ， 设 t=O 时刻点扰动源位于 0 点，在 t = 3 时刻 ， t=O 时的扰动到达半径为 3ct 的 0 3 球面上，而扰动点的水平移动距离为 3 凹，到达 M 点，由于 U < c ， M 点位于0;. 阑球 面内.在 t = 1 时刻，扰动点到达()' ，它的水平移动距离为 Ut ， 由该时刻发出的扰动在 t=3 时刻到达半径为 2ct 的回球面。"它位于图球面 ()3 内;同样道理，在 t=2 时刻， 扰动点到达 (j' ， 它的水平移动距离等于 2Ut ， 由它发出的扰动在 t = 3 时刻到达半径 为 ct 的圆球面 0 1 ，它包含于 ()2 、 ${\rm ο}$3 内.以上分析表明，亚声速运动的点扰动源，扰动 点始终位于扰动披内，在足够民时间以后，它的扰动总可以传播到整个空间.因此亚 声速运动的点扰动源的影响域也是全流场. (3) 超声速运动的点扰动的影响域(图7. 3(c)). 此时扰动点运动速度 U 大于声

U (a) (b) ${\rm •}$ (c) 图 7.3. 气流中扰动影响域

\[(8) U=O.Ma=O; <b) U<C.Ma<l; (c) U>C.Ma>l\]

速 c. 设 1=0 时刻点扰动位于 0 点，在 t=3 时刻该点扰动传播到半径为 3ct 的 03 球 面上，而扰动点的水平移动距离为 3 凹，由于 U>c ， t=3 时刻扰动点 M 位于球面 0 3 外.在 t=1 时刻，该扰动点到达 0' ， 它的水平移动距离为 Ut ， 从该时刻发出的扰动在 t=3 时刻到达半径为 2ct 的回球面 ()z， 它位于圆球面 01 的上游;同样道理，在 t=2 时刻，该扰动点到达 d' ， 它的水平移动距离等于 2 凹，由它发出的扰动在 t=3 时刻到 达半径为 ct 的阿球面。 l' 它位于 ()2 、 $\alpha$ 两个球面的上游.由于扰动披面的半径和声 速 u 成正比，受扰空间球心的位移和点扰动源的速度 U 成正比，因此，通过 M 点可以 作图矶、 O2 、。可、…的公切线，在空间中它是一个圆锥面，通过三角关系，不难求出锥 的半顶角 $\mu$ = arcsin(c/U) = arcsin( 1/ Ma) ， 称该回锥为马赫锥，半顶角 $\mu$ 为马赫角. 以上分析表明.以超声速运动的点扰动只能在下游马赫锥内传播，而不能传播到马赫 锥外.就是说以超声速运动的点扰动源的影响域在它的下游马赫锥内. 以上的分析很容易推广到绕流问题，如果在随点扰动源运动的坐标系中考察上 述扰动的传播，就可以观察到均匀气流(亚声速或超声速)受固定点扰动的流场情况. 这时影响域相对于扰动点的关系不变. 下面还可以通过影响域的概念来分析接收扰动的依赖域. 定义 7.3 空间固定点 P 能够接收到气流扰动信号的区域称为 P 点的依赖域. 从图 7.3(b) 中可以看到，在亚声速气流中任意点的依赖域是全流场，因为无论 是在指定点上游或下游的扰动，只要有足够长的时间，这些扰动总能传播到该点.而 超声速气流则完全不同(参见图 7.4) .P 点下游的扰动不能影响 P 点的气流状态;不 仅如此，上游的扰动只有一部分能够传播到 P 点.具体来说，以 P 点为顶点，以马赫 图7.4 超声速气流中的依赖域 z 倒马赫锥

角为半顶角，向扰动点上游作回锥(称为倒马赫锥) ，只有在倒马赫锥内的扰动才能传 播到 P 点，在倒马赫锥外的扰动，不论时间多长，都不可能传播到 P 点.因此 P 点的 依赖域是 P 点的倒马赫锥.图 7.4 中还表示了 一 点 P 的影响域. 总之，亚声速气流流场和不可压缩流场类似，扰动的影响域和依赖域是全流场. 超声速流场中扰动的影响域只限于扰动下游马赫锥内; 一 点的依赖域在倒马赫锥内. 这种力学特性对于分析和计算超声速流动是很重要的，例如，求解不可压缩或亚声速 绕流问题，必须给出绕流边界的全部边界条件，因为，从物理上说，上下游的条件都对 绕流有影响.而对于理想气体的超声速绕流问题，下游边界条件是不必要的，因为下 游边界流动参数的任何变化(扰动)都不影响上游的气体运动.

\section*{7.3  理想气体等摘流动的主要性质}\addcontentsline{toc}{section}{7.3 理想气体等摘流动的主要性质}

1.完全气体流动的基本性质和克罗科 (Croco) 定理 (1)理想气体定常绝热连续流动中沿流线恼不变 证明 在绝热、连续流动的条件下，理想气体质点沿迹线摘的变化率等于零 (式 (7.5 门，定常流动中，流线与迹线重合.因此，每一条流线上的蛐值是常数. (2) 理想气体绝热定常流动沿流线 h +U2 / 2 = const. . 证明 由能量方程 (7.4): ;)(h+U 2 / 2)/ r1 t+U. V(h+U 2

\[/2) = (;) \rho  / J t) /p+q , \]

在定常和绝热条件下，右端两项和左端第一项都等于零，因此有 U ${\rm •}$ V (h + U2 / 2) = 0 或 h + U2 / 2 = C($\eta$) (7.17) n 为流线参数.证毕.式 (7.1 7)是理想完全气体在流线上的伯努利积分. 定义 7.4 单位质量的蜡和动能之和 ， h ，， = h+U 2 /2 ， 称为总蜡 z 并用 h () 表示. 式 (7.17) 又可解释为理想气体绝热定常流动中沿流线总惜不变.在定常流动中， 总恰是同一流线上速度等于零处的馆值，故又称滞止婚. (3) 克鲁科定理及其应用 定理 7.1 理想气体的绝热定常流动中，若质量力可忽略不计，在全流场有下式 成立 z

\begin{equation*}.  XU = TVs-Vh o\tag{7.18}\end{equation*}

证明 由 Lamb 型方程出发 au , \_".. , $\tilde{I}$ U 2 \textbackslash\{\} 1 $\tau$::-+nxu+v 1-... )=一 $\tilde{Vp}$ rJ t 飞 " I P ;)u 定常流中万 =0; 根据热力学关系: Tds= $\rho$d (l / p) 十 de = 一 句/ p+dh. 或 - dp /$\rho$= Tds-dh. 因此在均质平衡气体中有一 Vp/p= TVs - Vh ，代入兰姆型方程可得

。 X U = TVs - V (h + 号) = TVs - Vh $\upsilon$ 证毕. 定义 7.S 辅值处处相等的流场称为均铺流场 z 总熔处处相等的流场称为均培 流场.对于均恼均馆的定常流动，由克罗科定理可得出以下推论. (1)轴对称或平面流动中均情均焰场必元旋. 均蝙流场中 Vs=O. 均始流场中啊 。 = 0. 因此根据克罗科定理，均铺均始流场中 应有。 xu=o. 在平面或轴对称流动中 n l.. U ， 从而 In XUI = Inllul =0. 流场中 U 子士。〈否则就是静止流场) ，因此必有$\beta$ = 0 ，即流场是无旋的. (2) 定常 三 维理想气体的均情无旋流动必均摘. (3) 定常 三 维理想气体的均恼元旋流动必均始. (4) 非均瑜的定常气维理想气体均饷流动必有旋. (5) 非均娟的定常 三 维理想气体均始流动必有旋. 推论 (2). 推论 (3) ，推论 (4). 推论 (5) 可由克罗科定理直接导出，不一一重复. 克罗科定理在理想气体的定常流动中有重要意义，利用它可以判断流场是否有 旋.只要流场中存在总始梯度或摘梯度，流动就有旋.又如，前面已证明:理想气体绝 热定常连续流动 '1' 7-荷流线总馆和愤不变.如果 二 维理想气体绕流的来流速度、总结和 情处处相等，则流场处处总培和辅相等，即全场是均熔均惰的.因此流动是无旋的.这 里流动连续是 一 项很重要条件，在高速气流巾，当气流速度超过当地声速时，可能产 生激波，气体质点穿过激波时热力学状态发往突变，是不可逆过程，这时虽然气体运 动是绝热的，但它的蛐将发生变化.因此定常超声速气流中有激波时，即使激波前流 场是均摘无旋的，激波后流场可能是非均铺有旋的. 2. 理想常比热容完全气体沿流线的等煽关系式 理想流体绝热定常连续流动时，沿流线铺值不变，称之为定常等情流.在很多实 际流动中，如 z 物体在静止大气中匀速运动;管道或叶轮机械通道中的气流运动等， 当流动的粘性效应可忽略，流场与外界无热交换又不出现激波时，这些流动都可视为 定常等情流动. (1)理想常比热容完全气体定常等摘流动的基本方程 定常绝热流动的能量关系式就是伯努利积分: 巴 +h = 巴 +$\iota$ T= 巴+ -y 主=巴 +iL = C(n)(7.l9) 俨 2 ' y - 1 P 2' y - ] 式中 c 是完全气体的当地声速.完全气体的定常理想绝热连续运动中的等情方程可 写作: 车 = B(n) P.

式中 C( 的 ， B(n) 是同一条流线上的积分常数. (2) 滞止参数、最大速度和临界参数 定义 7.6 在定常流动中，气体流动等情地减速到速度等于零的状态称为滞止 状态，滞止状态的气流参数称为滞止参数. 滞止参数记作户，如 ， To ， 儿，分别称作滞止密度，滞止压强，滞止温度和滞止熔. 它们可由等情流动能量公式和等恼过程确定.首先由能量方程 ho=cpTo = 【J2 /2+ h=U Z /2+ 马 T 导出滞止温度公式如下. ${\rm 1}$滞止温度 r rZ To = 芸一 +T C. C H 庐 其他滞止状态参数由滞止温度导出如下: \textasciitilde{} = (号)声，号= (号)击， So - S = $\iota$In(\textasciitilde{} )一 $\iota$In(; ) 由 To = 乒 +T 可以看出 z 滞止温度是理想气体沿定常流动流线的最高温度.

\[\iota  C p\]

(7.20)

\subsection*{例 7.1  在静止空气中以 800m/s 等速飞行的火箭，计算火箭头部驻点的温度和}

静止空气温度之差. 解 在等速运动火箭中的观察者，火箭周围空气是定常绕流，火箭头部的驻点是 滞止点.应用式 (7.20) ，火箭头部驻点温度和静止温度之差: To 一 T= 乒空气的比

\[t.C \rho \]

定压热容 c$\rho$ =yR/(y- l) = 1. 4 X 287/0.4= 1004. 5N ${\rm •}$ m/(kg ${\rm •}$ K) ，因此 T。 一 T = UZ /2c p = 318K 由以上计算可见，高速运动物体的头部有很高的温度，物体运动速度越快，头部 温度就越高. ${\rm 2}$理想气体定常等情流动中的最大速度. 定义 7.7 在理想常比热容完全气体定常流动中，流体质点等恼地加速到 h=O ， 或 T=O ， 或 $\rho$=0 时的速度为最大速度，记作 Umax ${\rm •}$ 由能量方程式 (7.19) ，令 h=O ， 得 Umax = /2瓦=/2$\zeta$T o = /2页王， /(y- l) ${\rm 3}$临界参数. 定义 7.8 在理想气体定常等情流动中，流体质点速度等于当地声速的状态称 为临界状态.临界状态下的气体状态参数称为临界参数，分别用 p' ,p' ,T' ， h' 来表 示临界密度，临界压强.临界温度，临界馆等. 根据定义:临界速度等于当地声速，即 U' = a' = .;百勺p* = .;页 T' = ..j(y-!)h' (7.21)

其他临界状态参数导出如下，由能量方程 : U. 2 /2+ 扩 =h o 和式 (7.2 1)导出

\[ho = (r+ l) h' /2\]

于是有 2一川 扩 -h一-

\[T-TH\]

(7.22) 由等情关系式得 主二 = f 主二 lfT=/ 一旦-IZZT po 飞 Tnl \textbackslash\{\}r+11 i\_ = f 主二沪 =/-L lh 向飞 To 1 飞 r+ 1I

\subsection*{例 7.2  双原子气体的临界参数值.}

解 双原子气体的绝热指数 $\gamma$= 1. 4. 因此 主二 f\_2\_ \textbackslash\{\}f"i = (l <;?R'L T' =1 一一--:- I = O. 5283. 一一= 0.833. 一= 0.634 Po \textbackslash\{\}r+11 -.----, To -.---, po (3) 用马赫数表示的完全气体等'庸关系式 气体定常等摘流中流动参数的变化可以用能量方程导出，对于完全气体，这些公 式还可以用马赫数作为自变量来表示，它们更便于分析和计算.推导过程如下 z 由滞 (7. 23) (7. 24) 止温度公式 Tn =U 2 /2c p +T. 可导出 工= 11+ 卫二) -( T o 飞 2c p T 1 对于完全气体 .cp=rR /( $\gamma$-1) .因此 U 2 /(2 句 T)=( $\gamma$ -1)Ma 2 /2. 代入滞止温度公 式，得 T h I. ， Y-] 八 l 于;=瓦= (1+亏二地 Z ) 进一步代入其他状态公式，就有 (7.25) ;=(1+ 号l Ma 2 ) f${\rm i}$ (7.26) 去= (1 +平叫声 (7.2 7) 去= (1 +早地 2 r+ (7.28) (4) 气体等情流动的不可压缩近似 一般来说，气体是可压缩的.下面将说明，当理想气体定常流动的马赫数很小时， 气体流动可近似为理想不可压缩流动.首先导出以下压强系数公式 z

\[\rho  o ,\]

4二主=王二=丁$\iota$1(1 +乓l Ma 2 户一 Il $\div$pI 尸 号 ylV1 a - L 、 $\iota$ ， J 式中 U 、 $\rho$ 、 Ma 是气流的当地速度、压强和马赫数 "${\rm þ}$u 是等情滞止压强.当 Ma<1 时. 我们可以将 (1 +早w 广作泰勒级数展开，得(由读者自己完成计算) $\rho$。 一 $\rho$lT一一!:.=1+--:- 岛企1 2 + ()(岛也 I ) 言p[P 啥 当 Ma<<1 时，上式可近似为: (p。 一 $\rho$ ) / (p 1Y /2) = 1. 即

\[p" = \rho  +plP / 2\]

这就是不可压缩理想流体定常流动的能量积分，或不可压缩理想流体的伯努利积分， 根据式 (7. 29) 可以估计适用不可压缩近似的马赫数范围.例如 z 当 Ma = 0.3 时，近 似公式所得结果与伯努利公式之差为 0.0225. 即 2 \% 左右，这对工程应用是允许的. 因此，当流场中最大马赫数不超过 0.3 时，理想气体的定常等摘流动往往可以用不可 压缩理想流体方程来近似. (5) 速度系数 定义 7.9 流体速度与临界速度(或临界声速)之比称为速度系数，记作 (7. 29)

\[u-f\]

一一 A 速度系数 A 与马赫数 Ma 之间有一 一 对应关系.速度系数可写作 ${\rm A}$ = 严=?去? = 存在于 将前面已导向的关系式 T' /T，， =2/(y 十 1 )和 T/T，， =[I+(y 一 1 )Ma 2 /2J I 代入上 式，得 $\lambda$= 协[击 (1 +平Ma 2 ) ] ' (7.30a) 或 「 $\gamma$ + 1/ ， $\gamma$- 1, 2 \textbackslash\{\} l $\tau$ 儿也 = ${\rm A}$I 一一$\tilde{r}$ 1 一一一-二 ${\rm A}$ " 1 1 L 2$\gamma$y 十 1" J J 速度系数 A 与马赫敬之间的关系示于图 7.5. 通过计算导数 d$\lambda$ (Ma) / dMa 和 dMa (，\textbackslash\{\})/ 心，可以验证 $\lambda$ ， (Ma)>O 和 Ma'( $\lambda$)> 0 ，即 HMa) 和 Ma(${\rm A}$) 都是正值的单调增函数，如图 7.5 所示.而且还有以下关系 z Ma=O 时 .${\rm A}$=O;Ma=1 时 .${\rm A}$=I;Ma>l 时 .${\rm A}$>I;Ma<1 时 .${\rm A}$<1 以及 (7.30b) 地$\rightarrow$ $\infty$ 阳$\rightarrow$ J吉

\section*{7.4  激波理论及应用}\addcontentsline{toc}{section}{7.4 激波理论及应用}

2.5 2 1.5 ，哇 0.5 Mu 图7.5 速度系数与马赫数之间的关系 <y= 1.的 以上关系说明，我们也可以用速度系数来判断超声速()，>1)或亚声速 ($\lambda$< 1)流动.对 于常用双原子气体 r= l. 4 ，这时 $\lambda$( $\infty$)=./6. (6) 用速度系数表示的等情关系式 由速度系数和马赫数间的关系，还可有以下的等情关系式 z T\_/..!:...\textbackslash\{\}2=I土.!f l -I二.!12 \textbackslash\{\} T 盹飞 c' J 2$\gamma$r+ 1" , (7.3 1) ..l... 丘= iI土.!(1- 1二\textasciitilde{}${\rm A}$2) l Ti (7.32) p' L 2 飞 r+ 1" J J \textasciitilde{}= i I土.!(1- 1二与2 门 f${\rm i}$ (7.33) $\rho$ 到 L 2 飞 r+ 1"' J J 为了实用方便，完全气体等恼流关系式已制成表楠，给定马赫数后，其他参量 ${\rm A}$. $\rho$/ 扣 ， p/户 ， T/T，。 可以从表中查出.见本书附表 4.

\section*{7.2  节讨论过气体中微弱扰动披的传播过程.在很多实际问题中，例如:飞行器}\addcontentsline{toc}{section}{7.2 节讨论过气体中微弱扰动披的传播过程.在很多实际问题中，例如:飞行器}

高速飞行、超声速气流的突然压缩等流动中，常常出现强扰动，这时往往会出现流动 参数的突变现象 z 激波.在这一节中，我们先从物理上简要说明激波的形成过程，然 后建立激波前后物理量之间的关系式. 1.正激波形成的物理过程 考虑一种简单的气流压缩过程.在无限长的等截面管道中充满静止的可压缩气 体，初始状态为 : u=O.p= $\rho$ " ， p= 酌 ， T=T". 管的左端有一活塞从静止开始向右运 动，速度逐渐加大到 U( 不是小量) .然后以等速 U 运动.设想活塞加速过程较慢，与

活塞表面靠近的 气 体不断引起微弱扰动，这些微弱扰动以 声 波的速度 一 个个向右传播. 第 一 个波以 C o 速度向右传播(见图 7. 6) .并使波后气体产生 一 个微小的速度 UI 图7.6 说明激波产生用阳 (等于当时当地活塞速度) .而且因气体受到压 缩，温度、压强和密度都略有增加，波后受压缩 气体的声速 ('1 大于原静止的声速 c u ${\rm •}$ 活塞继续 加速而引起第 二 个小扰动声波 H 才，它在以速度 U I 运动的气体小以声速向向右传播，故第 二 道 波的绝对速度是 ('1 + UI > C ，，; 第 二 道波过后，流 体继续向有加速为 U 2 ( > UI ) .同时温度 Tz > T ， ${\rm •}$ 因此后续波 (' 2 + U Z> CI + UI > C " . 依次类推， 当活塞不断向右加速时;第四道波，第五道 波，...... .一道接 一 道的扰动波向右传播，而且后续扰动波的波速总是大于先行波的 波速，所以后面的波 一 定会追上前面的波.经过 一 定时间后，无数个小扰动弱波叠加 在 一 起形成 一 个有限强度的扰动波 一一 激波，它将以大于静止气体中声波的速度向 右传播. 2. 驻正激波前后物理量间的关系式 (1)激波的简化模型 上面定性地说明了理想气体中激波生成过程.实际上激波是很薄的 一 层，它的厚 度是分子自由程的量级.在这 一 薄层巾，物理量(速度、温度、压强)迅速地从波前值变 化到波后值，速度梯度、压强梯度和温度梯度都很大，因此研究激波层内流动时必须 考虑粘性和热传导的作用.由于激波厚度相对于流体的宏观运动非常之薄，在实际问 题中，我们对激波内流动情况并不感兴趣.而只需知道物理最通过激波后的变化就足 够了.因此我们忽略激波厚度而将激波简化成数学上的间断面.由于激波层很薄，当 激波层中不发生离解、电离等物理、化学过程时 ， 气体穿过激波可以认为是绝热过程. 就是说，可以将激波简化为绝热的间断面. (2) 正激波的相容条件 定义 7.10 和气流速度垂直的物理量间断面称为正激波. 一般情况下，激波可以在气体中运动，也可能是驻定的，为了便于分析，始终将坐 标系固结在激波上，将正激波看成是静止的平面，这种激波称为驻激波.这时气流穿 过激波，气流参数在激波面两侧发生跳跃.由于物理量在激波面上间断，不能用微分 形式基本方程来分析这种现象，只能应用积分形式的基本方程来建立激波前后的物 理关系式. 取 一 包含激波面的很薄的长方形控制体(垂直于图面方向取单位长度) (如 图 7.7) .控制体侧面积 A ， = 儿，且平行于激波面 .A ， 和 Az 间距离 d .r$\rightarrow$ O. 是上下控制

\section*{7.4  激法理论及应用}\addcontentsline{toc}{section}{7.4 激法理论及应用}

面的宽度.在控制体上应用守恒定理，且当 dx $\rightarrow$ 0 时， 控制体内质量、动量和能量的体积分值为小量可忽略， 因而激波前后有以下关系式. 质量守恒方程

\[P1 U1 = pzUz\]

动量方程 $\rho$ l+p1m = 扣 +pzm 能量守恒方程 T 12 T 12 (7.34) (7.35) (7. 36a) 激 波 什ti 2 屿 I I I T2 T飞 c 凡+亏 = c${\rm þ}$T Z + 亏 图 7.7 推导激波关系式用图 或 状态方程 r P1 L U\textasciitilde{} \_ r pz \_j\_ m 一二一一十一

\[r-1p1' 2 \gamma - 1 pz , 2\]

扣一 $\rho$I RTzpz RT 1p1 (7. 36b) (7.37) 以上 4 个方程是联系激波前后压强、密度、温度和速度的 8 个流动参数的基本方 程，也称为正激波的相容条件.应用以上方程可以分析激波前后气流参数间的关系， 一般来说，只要知道激波前后 4 个参数，就可以用以上关系式确定其他 4 个气流参 数.下面利用以上关系式分析激波的主要特性，并导出便于计算的激波前后气流参数 的关系式. (3) 激波过程的兰金-于戈尼奥 (Rankin e- H ugoniot) 关系式 由正激波的质量守恒方程 P1U1 =pzUz 可导出 U 1 \_ pz Uz P1 将它代人动量方程 (7.35) 得

\[\rho 1 - pz =.. P1 U1 (U Z - U, )\]

还可由质量守恒方程导出 U1 + U2 = U1 ( 1 +在)= P1 U, ( 卢+卢) = pzUz ( 卢+卢) 从而得 1 , 1 1 一+一 = \textasciitilde{}(U ， +U2 ) p, pz P1 U , 将动量方程 $\rho$ ，- pz =P1U1 (U 2 -U ，) 与上式相乘得 11 , 1\textbackslash\{\} ${\rm •••}$ ${\rm •••}$ (p ， 一 $\rho$z) I 一+-=-)= Ui -U\textasciitilde{} 飞 P1 pz I

再代人能量方程 (7.36b) 得 2)' I $\rho$I \_ pz \textbackslash\{\} \_ I \textasciitilde{} \textasciitilde{}， I 1 \_l\_ 1 \textbackslash\{\} 丁L$\tau$l 一一 y< 1= ($\rho$I - Pz)(\textasciitilde{}+\textasciitilde{} 1 1- 1 飞 Pl pz I 飞 Pl pz I 等式两边乘 PZ/PI 并进行整理可得 ()'+ 1) 旦 -($\gamma$- 1) 扣 PI PI (7. 38a) ($\gamma$+ 1)一句一 1) 色 PI 或 pz 一主立 +()'-1) PI ()'一1)生 +($\gamma$+ 1) 1' 1 由状态方程 (7.37) 可导出激波前后的温度比 z (7.38b) ( y + 1) pz + (y - 1) (主!. ) z Tz \_ PZPI 一， $\rho$l' 飞 PI /

\[T1 \rho  I pz (y + 1) pz + (y - 1)\]

PI 式 (7.38a) ，式 (7.38b) ，式 (7. 38c) 称为兰金-于戈尼奥关系式，也叫做激波绝热曲线， 它是气体通过绝热激波的过程方程. (7. 38c) (4) 激波过程和等蜻过程的比较 为了准确地了解激波，我们把激波绝热过程方程 (7. 38a) 、方程 (7.38c) 和等摘过 程(即连续绝热过程〉的热力学参数变化作一比较.在图 7.8 中同时给出了激波绝热 曲线和等楠过程曲线. 5 J 反， L 且 1. '3 '3 、，\&匀，${\rm ‘}$'3l

\begin{gather*}
a\cdot \\
/ 2 \rho \iota 
\end{gather*}

l ，、 d nu nu nu 5

\section*{4.5  ←一一 ←一-}\addcontentsline{toc}{section}{4.5 ←一一 ←一-}

11 I $\leftarrow$一一一- f v J 4 3.5 号 3 龟. 2.5 2 1.5 0.5 0 o 0.5 1 1.5 2 2.5 3

\[T/T1\]

(a) (b) 图 7.8 激波压缩和等精压缩的比较(实线 z 激波过程 z 虚线 z 等精过程〉 ( a) 压强比和密度比间的关系; (b) 压强比和温度比间的关系 由式 (7.38) 和图 7.8 可以得出以下结论. ${\rm 1}$激波压缩是有限压缩，气流通过正激波后，压强和密度都增高.但是，正激波

\section*{7.4  激波理论及应用}\addcontentsline{toc}{section}{7.4 激波理论及应用}

后的密度增高有极限，由式 (7.38b) 可以计算，当扣 Ipl $\rightarrow$$\infty$时， 1: 一 pz \_ r士11 \_" …一一二 - 2『;PIY-1| 卢 M V 而等铺压缩是无限的，当 $\rho$ 21 Pl $\rightarrow$$\infty$ 时 ， p2lpl $\rightarrow$$\infty$. ${\rm 2}$对于相同的压强比，激波压缩的温升比大于等摘压缩的温升比(图 7.8(b)). ${\rm 3}$从图 7.8(a) 可以明显地看到:激波绝热曲线和等摘曲线在 P2 年1=1 点相切. 这表明微弱的激波压缩接近等摘压缩波. ${\rm 4}$从图 7.8(a) 可以明显地看到\textasciitilde{} :相同密度比下 ($\rho$21pl >1)，激波压缩过程的压 强比大于等精过程的压强比;相反，假如出现激波膨胀过程时 (pzl Pl <1)，激波后与 激波前的压强比小于等摘过程的压强比.下面就要说明，激波膨胀过程是不可能出 现的. ${\rm 5}$激波压缩过程蜻增必大于零，是绝热不可逆过程. 由热力学公式，摘增为 A 叫 P 21PI \textbackslash\{\}一 $\omega$ 1n 「 Pz/$\rho$1 l 一、 lnl 扣 /$\rho$$\tilde{l}$ S = 52 - 51 二 CVl lpl/pff tvl L(户 1 Pl )y J = Cy In L 百万五$\tau$ 」 式中 \{$\rho$ 21 $\rho$1 ), = (pz 1 Pl )y表示等情过程的压强比.由图 7.8(b) 可见，激波压缩时 扣 /$\rho$1> 1，根据性质${\rm 4}$激波绝热曲线上 P 21 PI > (户 21 $\rho$l$\lambda$ ，将它代人铺增公式，可以 作出结论 ， S2- S I>O. 即激波压缩是情增过程，因此是绝热不可逆过程. @激波膨胀是不可能的. 如果激波过程中出现 pzlpl<l ， 激波后的压强小于激波前压强 ， P2lpl <1 ，即出 现激波气体膨胀的情况.由性质${\rm 4}$可知如 1 PI <(P2 年 I ), ，因而 52- S 1<O ， 就是说绝 热膨胀间断是情减过程.在热力学中我们知道绝热的情减过程将导致荒谬的第二类 永动机，这是不可能的. 因此可以断言:激波只能是压缩波，不存在膨胀激波. ${\rm 7}$激波前的速度必大于激波后的速度. 由于激波是压缩过程，激波后的压强和密度必大于激波前的压强和密度 ， pzl p ， > l ，于是由质量守恒关系式就可以导出 U2 I U 1 <1 ， 即激波前的速度必大于激波后 的速度.总之 z 气流穿过激波时气体压缩并减速. (5) 普朗特关系式 将动量方程 (7.35) 除以连续方程式 (7.34) ，并考虑 C 2 =yplp 可得 4+U1= 兰 +U2 y$\upsilon$ I yU 2 应用临界参数的定义，可由式 (7.36b) 得 d=L土14$\cdot $2 - Z二lu\textasciitilde{}

\[1 - 2' 2 '-'1\]

d=L土1c 叫 -L二lu\textasciitilde{}, 2 - 2-'

将以上两式代入前面的公式可得 t+UI= 云 +Uz 即 (U, -U?)(l-.\textasciitilde{} 二 1= 1 -.飞 U1UZ J 由于有眼强度的激波前后 U1 =I= 叭，故有

\begin{equation*}U1U2 = (".2\tag{7.39}\end{equation*}

这就是说正激波前后速度的乘积等于临界声速的平方.根据速度系数的定义，可立即 得到

\begin{equation*}\AA I\AA Z = 1\tag{7.40}\end{equation*}

这就是普朗特关系式. 因为激波前速度大于激波后速度，正激波前后必有 ${\rm A}$I>${\rm A}$Z' 于是由式 (7.40) 可得 出结论:正激波前为 $\lambda$.> 1<即 Ma.>l) 的超声速气流，正激波后为 $\lambda$z<H 即 Maz< 1)的亚音速气流.换言之，只有超声速气流才有可能产生激波. 综合以上分析结果，激波有以下主要性质. ${\rm 1}$激波是:幅增压缩过程. ${\rm 2}$只在超声速气流中才能出现激波. ${\rm 3}$激波是减速过程，驻定的正激波后是亚声速气流. 理解激波过程的定性分析后，利用以上公式来导出激波前后气流参数的关系式. (6) 正激波前后流动参数间的关系式 ${\rm 1}$激波前后的速度关系式和 Ma 数关系式 将 A 与马赫数的关系式 (7. 30) 代入式 (7. 40) 可得激波前后马赫数之间的关 系式 z 激波前后速度比 y-1 1+ 一$\tau$一地 i Ma\textasciitilde{} 一』

\[dyM Z Y-I a j -\]

2 一 (7.4 1) UI - ui - ui -f= (y+ 1)品也 i 2 (7.42) 瓦 - U1UZ -? -AI - 2+(y- 1) 地 1 ${\rm 2}$激波前后压强、温度等参数间的关系式 激波前后密度比由连续方程 (7.34) 导出 z 向 U， 一 (y+ 1)必1a\textasciitilde{} Pl U2 2 + ($\gamma$- 1)岛也 i (7.43) 激波前后压强关系可由动量方程 (7.35) 导出 z

主二主!.=巳旦 11 一旦 l = y$\tilde{i}$ 11 一旦丁 = yMai 11 一旦 1

\[\rho  ,  p, L\tilde{U} ,J a\tilde{L}' U ,J "'--'L- U ,J\]

将式 (7.42) 代人并整理得 生= 1 + 2,Y. (Ma \textasciitilde{} - 1) $\rho$1 Y 十 1 '-'\_, (7.44) 最后，激波前后温度比由状态方程求出 z T 2 \_ P2 Pl 一[ 2yMai 一 ($\gamma$ 1) J[($\gamma$-1) 儿也 i + 2J T, $\rho$ 1 P2 ($\gamma$+ 1) 2 1\textbackslash\{\}也? (7.45) ${\rm 3}$激波前后的滞止参数及性质 由激波前后的能量守恒关系式 (7.36a) 可知，激波前后的总惜不变，因此滞止温 度相等，即有

\begin{equation*}T 02 = T o,\tag{7.46}\end{equation*}

这就是说激波前后滞止温度相等. 滞止压强关系可由坠王 =P02 ( Ma 2) 生 (Ma ， )主!..(Ma ， )计算.这里应当指出，气流

\[\rho 01\rho 2\rho 1 . POI\]

穿过激波面时是情增过程，但在激波以前和激波以后的流动都是等惰的.因此激波前 的滞止压强比如 Ilpl 和激波后的滞止压强比 $\rho$ 021 P2 都可用前面给出的等摘关系 式 (7.27) 来计算.而扣/户 l 应当用激波关系式 (7.44) 代入，经整理可得 I y+1u二， 2 l 卢T $\rho$02 一 1\_2\_. 一\_' I I \_\_1y\_ \_ J\textbackslash\{\} A\textasciitilde{} 2 \_ Y - 1 l 元lT1 - 1 1 1 一一地 一一一 l' (7.47) 如 11 +宁叫 Ly 十广 l 川」 由简单的代数计算可得 ， Ma ， = 1 ,P021 POl = 1. 当 Mal>l 时， $\rho$021 Po, <1 ，就是说通过 激波后滞止压强减小(参见图7.9). 激波前后滞止密度可通过状态方程得到，因为 激波前后滞止温度相等，所以激波前后的滞止密度 比等于滞止压强比，即 色王=主旦 (7. 48) 户JI$\rho$01 以上关系式 (7.37)\textasciitilde{} 式 (7.48) 均已做成函数表 (见本书附表 5) 以便 t 程计算，只要知道 Ma ， ,Maz' $\rho$ ， P2 P2 U , T 2 P 一，一，一-一，」中任意一个，便可查得激波$\rho$oz pIPI'UJ TI 户。 l 前后其他流动参数之比. (7) 应用举例 图7.9 激波前后的滞止压强比 随马赫数的变化

\subsection*{例 7.3  管内超音速气流速度测量如图 7.10 所示，在超声速气流中放置一钝头}

的测速管，这时在测速管前产生正激波，已在测速管前的管壁上测得 $\rho$ ，=

\[10. 5kN/m\]

2 ${\rm •}$ 激波后气流的温度 T2 = 453. 3K. 以及测速管头部测得激波后总压

\[\rho 02 =50kN/m\]

2 ${\rm •}$ 求 Ma1 ， U1. 解 由给定参数得 主!... = 0.21

\[\rho 02\]

查附表 5 得 : Ma1= 1. 82 ， T2/T1= 1. 547. 故得

\[T1 = 453.3/1.547 = 293(K)\]

CI = v'页"T;" = 343. 1m/ 日 ${\rm •}$ U 1 = C1Ma 1 = 624. 4m/s 灿 1>1 e\textasciitilde{} PI P02 r f 图 7.10 管内超音速气流速度测量示意图 U 1 图7.11 兴Uj \_T01 $\rightarrow$/气!\textbackslash\{\}!?1 超音速进气道示意图

\subsection*{例 7.4  超音速飞机在高空等速飞行，空气经正激波(相对于飞机不动)进入超}

音速扩压器(如图 7.11) .将坐标固结在正激波上(即在飞机上) .测得 U2 =260m/s. T02 =400K. 认为正激波后流入超音速扩压器的气流速度 .T，$\omega$ 为正激波后气流的滞 止温度. 试求: (1)在此坐标系内的 $\rho$"2/$\rho$01; (2)U 1; (3) 若 T02 不变，问 Uz 增加到多大 时，扩压器前不出现正激波? 解 (1)已知空气的比定压热容 c$\rho$= 1006NM/( kg ${\rm •}$ K) ，由 TU 2 = T z +U\textasciitilde{} /2 c $\rho$ 得 T 2 =366. 4K ,C2 = v'页 T2 =384m/s. 由矶和 C2 算出 Ma2 =U2/C2 =0.677 1. (2) 由正激波表(附表 5) 可得 儿也 1 = 1. 57. T\textasciitilde{}/T ， = 1. 367. $\rho$02/ $\rho$川 = 0.9061 然后可计算出 z 2 ${\rm Æ}$ = 328. 4m/s. U 1 = CI 比 1 = 515 叫s (3) 当 Ma ，运 l 时不由现正激波，临界值为 Ma=l. 此时有 E主毛::; =斗(1 + 与1斗) 1 利用能量方程得: U z = ';2c $\rho$ ( T02 一 T2 ) =366. 24m/s. 即 U2 >366.24m/s 时，发动机 前未出现激波.

3. 运动激波及反射 实际中常有激波在静止空气中运动的情况，例如强爆炸产生的冲击波在空中的 传播等.采用惯性系中力学和热力学定律不变的原则，选择 一 个等速运动的惯性系， 使激波在该惯性系中是驻定的，然后利用驻激波面前 后的流动参数关系式导出运动激波前后流动参数的 关系式.例如当激波相对于静止坐标系以某 一 速度 D 运动时，只 需 取以 D 速度运动的坐标系，使激波相对 于新坐标系静止，然后应用驻激波前后 气 流参数间的 关系进行计算.以图7. 12 为例来说明如何计算运动 激波前后的流动参数. 酬〕一动一

\[U'2=D-U2\]

L " $\tilde{t}$ 激波 图 7.12 运动激波及坐标变换

\[U'I =[J-U1\]

当激波相对于静止坐标系 .:r 以速度 D 向左运动 时，激波前、后气流速度分别为队 ， U2 ${\rm •}$ 倘若在以激波速度相同的运动坐标系 z' 中来 观察和分析，这时激波是驻定的，驻激波前后气流相对驻激波的速度分别为 U\textasciitilde{} 和 U\textasciitilde{} , 可以算出 U\textasciitilde{} = D - U , ， U\textasciitilde{} =D - U2 ( 当波前绝对速度 U1=O 时 ， U\textasciitilde{} =D); 在运动坐标 系中，除滞止参数外其他热力学参量 和驻激波相同.具体计算如下. (1)静止气体 rt 1 运动激波速度和伴随速度 (U 1 = O ， U\textasciitilde{} = D) 假定运动激波前气体是静止的，气体状态是 $\rho$ " p ， .T ， .已知激波前后的压强比 为如 / 扣，下面来计算激波运动速度和激波后的 气 流速度. 首先，可以从驻激波前后的压强公式导出运动激波前的气流速度(驻激波的速度 和马赫数都用上标撇号表示) .由式 (7. 44) 生 = l+lL( 胁 '12 \_ 1) $\rho$1 -, Y + 1 , - - ${\rm •}$ 可得 儿也" = .IY \textasciitilde{} 1 + Y \textasciitilde{} 1 $\rho$ z 一'一于一- a 叫V 2y , 2y $\rho$I 于是 - \_, /y 一 1 , y+ 1 扣D = U\textasciitilde{} = (J\textbackslash\{\}.企1'1)('1 = ('1./ 一一一十一一一一 勾 2$\gamma$ 2y PI 驻激波后的速度 U; = U:p， 1户，在固定坐标系中，激波后的速度由运动转换关系求出 z

\[U z = D - U~ = U~ - U~ = U~ (1 - PI 1 P2 )\]

将激波前后密度比公式 . l!3\_ = $\tilde{Y}$ 十 1 )$\rho$ 21 $\rho$1 + <y 一 1) 和激波运动速度 D 的公式

\[PI (\gamma  - 1) p zl \rho 1 + (y+ 1)\]

代入，得运动激波后的速度为 $\iota$ = (.', ${\rm •}$ /主 $\rho$ 2 1 $\rho$1 一 1 " V Y .j (y- 1) +(y+ 1) 扣 / 户 l

在静止气体中传播的激波后气流速度称为激波伴随速度，或简称伴随速度.当强激波 通过时，激波伴随速度可高达每秒几百米，它也是一种相当大的爆炸破坏力.伴随速 度也可用驻激波前马赫数和声速表示.将向 /pJ =[($\gamma$ + l) Man/[2+(y 一l) Man 代 人激波运动速度 D 的公式，可得 u. = U\textasciitilde{}-U\textasciitilde{}= \textasciitilde{} \textasciitilde{}(地 :2 - 1) l. - I - L r 十 1 必知 ll - 1JIJ7.S 设爆炸前后压强比等于 10. 激波前是常温大气 (T=293 仙，计算激波后 的伴随速度. 解 常温下的气体声速 CJ = VyR(273+20) =343m/s. 代人伴随速度公式，得 U2 句 747m/s (2) 行进正激波的平面反射 在静止空气中一道行进激波，当它遇到障碍物时将从障碍物上反射.例如一道平 面激波遇到一堵直墙，当激波撞到直墙时，激波后的伴随速度在墙面上为零(因墙固 定不动) .这时必有激波从墙面上反射出来.下面，利用以上公式计算激波反射过程.

\subsection*{例 7.6  已知激波在壁面反射前的波前参数为 zρJ = 100kN/ 时 • T J = 287K. 波}

后压强为扣 =450kN/m 2 ${\rm •}$ 计算激波在壁面反射后被后压强、温度和速度(如 ${\rm •}$ T3 参见 图 7. 13(b)). I PJ P 2 .T1 ， U1 卡- TJ I .叫 UJ=O 人射撒波 绝热固壁 p 、 p、 .T， .U、--$\rightarrow$寸 T，

\[D1 I U,=O\]

反射激波 (a) (b) 图7.13 运动激波及其.反射 (，，)人身.t 波; (1))反射波 解先算人射激波前后参数. f色!.\textasciitilde{}凶|瞠 (1)根据给定参数和运动坐标系 ， cJ=1页\textasciitilde{}=340m/s. 激波前后的压强比 W ${\rm þ}$J =4.5. 由正激波表(附表 5) 查得驻激波前后参数为 Ma\textasciitilde{} = 2. 协\textasciitilde{} = 0.5774. T 2 /TJ = 1. 688 因此在固结于激波上的相对坐标系中 : U'J = Ma$\tilde{cJ}$ = 680m/s. T 2 = 485. 6K. C2 = J页T2 =44 1. 7m/s.U\textasciitilde{} =Ma$\tilde{c}$2 =0. 5774 JyRT 2 =255m/s. 入射运动激波速度和激波后气流速度(伴随速度)分别为 D J = U\textasciitilde{} = 680m/s. U2 = U\textasciitilde{} 一 U\textasciitilde{} = 425m/s

再算反射后的气流参数. (2) 反射波(向左运动速度 Dz) 后流场. 驻波前、波后气流速度仍用町 .U\textasciitilde{} 表示，这时 ， T\textasciitilde{}= 丑 ， T\textasciitilde{}=T3 ; 相应地 c;=cz ， c;=C3. 由静止固壁条件可得:反射波 Dz 后气流绝对速度 V3=0 ， 故在相对于激波面 (运动速度为 Dz) 静止的坐标系中，波后气流相对速度 U\textasciitilde{}= 鸟，波前气流相对速度 U; =Uz+Dz. 可以通过反射波前后速度变化求出对应于驻激波的马赫数和热力学 参数. 已知: U; - U\textasciitilde{} = au = 425m/s. 应用前面导出的 U;-U\textasciitilde{} 公式，可解出 Ma:= 1. 73. 再用正激波关系式(查正激波表)可得激波前后的参数比 z

\[\rho 3/\rho  z=3.325 ,  T3 /Tz = 1. 48\]

P3 一 $\rho$ :!${\rm Þ}$ 一最后得 z 一一 jJ:! I' Z =4. 5 X 3.325= 14.96. 瓦 =485. 6X l. 48=719K.

\[\rho l\rho Z \theta l\]

从上例的计算结果可以看出:激波在固壁反射后，压强急剧增加.本例中，运动 激波的马赫数仅等于 2 ，激波在静止固壁反射后的压强等于反射前的 14.96 倍;同时 温度上升很高.这说明激波具有极大破坏力. 4. 斜激波理论 在实际的超声速流动中，很多激波间断面与气流方向不垂直，称与气流方向不垂 直的平面激波为斜激波. 下面研究驻斜激波的情况，即斜激波间断面是固定的，激波前后速度是均匀的. 超声速气流绕尖模的流动(图 7.14(a)) 和超声速气流绕有小折角的平壁面的流动 (图7. 14(b)) 都能产生驻斜激波. 斜激波前后气流速度之间的几何关系示于图 7. 14(c) ，图中 U 1 ${\rm •}$ 、 Uz ${\rm •}$ 分别为激波 前与激波后气流速度在激波面法向的分量，矶，、 U2 ， 为激波前、后气流速度在波面切 向的分量，即平行激波面的速度分量.激波后气流方向与激波前来流方向的夹角 $\delta$ 称 为气流转角 z 激波间断面与波前气流方向的夹角 p 叫做激波倾斜角，简称激波角.当 卢=$\pi$/2 时，来流和激波间断面垂直，相当于正激波 i${\rm ß}$=$\mu$ 是微弱扰动的马赫角，所以 斜激波倾斜角 $\mu$ 〈卢〈 $\pi$/2. (1)斜激波相容条件(基本方程) 与导出正激波关系式的方法相同，取垂直于激波面的控制体，并在控制体上建立 质量、动量和能量方程. 质量守恒方程 $\rho$IU 1 ${\rm •}$ = pzUz${\rm •}$ (7.49) 动量方程 间断面法向的动量平衡式

\begin{equation*}\rho 1 +p1Ui. = \rho  2 +pzm.\tag{7.50a}\end{equation*}

\[Ma,>1\]

,/ ...../二'/ (a) (b) 图7.14 斜激波前后气流的几何关系 间断面切向的动量平衡式

\[p , Uln U" = p2 U2. U21\]

应用质量守恒方程 (7.49) ，切向动量方程可简化为 第 7 章 气体动力学基础 p (c) (7.50b)

\begin{equation*}U" = U21 = U,\tag{7.50c}\end{equation*}

这说明气流通过斜激波时，切向速度分量没有变化，或者说，切向速度是连续的. 能量方程 (7.51a) 盯? 口子 UL , U\textasciitilde{} ， -\_:" + -\_" +h , = -:旦+ -n u

\[+ h2 = h, 2 ' 2 ,--. 2 ' 2 '\]

由于 U 11 =U2 " 能量方程可简化为 r r2 r r2 r r2 于十 h ， - 子 + h2 = 11 0 一亏 - h\textasciitilde{} (7.51b) 或

\[(7.51c) y\rho  ,  , Uf" - r Pz , m ..\]

Y 一 1 p, '2 r - 1 P2 ' 2 状态方程 p, \_ P2 Pl T , pz T 2 (7. 52) 以上式 (7.49)- 式 (7. 52) 是气体通过斜激波的基本方程，下面利用这组方程分析斜 激波前后气体运动的性质，井导出计算公式. (2) 斜激波前后气流参数间关系式 将式 (7.49)- 式 (7. 52) 与正激波基本方程比较可看出 z 气流通过斜激波相当于 以法向分速度通过正激波的流动.只要将正激波关系式中下标"\}"，， "2" 换成 "ln" ， u: "2n" ，总惜 ho 换成 h\textasciitilde{} =11 一一，则可将正激波关系式转化为斜激波关系式.和推导。 2 正激波前后气流参数关系式方法相同，不难证明斜激波前后压强比和密度比的兰 金-于戈尼奥关系式与正激波是相同的.因此斜激波具有和正激波相同的基本性 质 z 斜激波是绝热摘增过程;斜激波后的压强和密度必大于斜激波前的压强和密 度 z 驻斜激波后的气流法向速度必小于驻斜激波前的气流法向速度.从运动学角度

来看，由于斜激波前后的切向速度相等，我们可以建立一个运动坐标，它以斜激波 的切向速度运动，在这个运动坐标系中气流穿过激波的运动和正激波相同.因此可 以利用正激波的公式推导斜激波前后的 气 流参数关系式，斜激波前气流垂直于间断 面的速度分量 等于 U1 ${\rm •}$ =U1 sin${\rm ß}$ ，它相当于正激波前的气流速度，相当于正激波前的 气流马赫数应为 U1./CI =MaISin${\rm ß}$. 利用正激波关系式，就可以得到用马赫数表示的 斜激波关系式为 fl2一 (r+ 1)具备2$\tilde{sin}$ 2 ${\rm ß}$ Pl 2 + (r- 1) 儿1a $\tilde{sin}$ 2${\rm ß}$ P2 = \_\_1I\_Mn 2 "in 2 f.1\_ r二1Pl r + 1" …,..…,. r+ 1 T2 一 r 2 讪${\rm ð}$a $\tilde{sin}$ 2 ${\rm ß}$ 一 (r - l ) l r2 + (r-l) 儿也$\tilde{s}$ in 2 $\beta$1 于\textasciitilde{}- L r+ 1 JL-($\gamma$+ 1)险$\tilde{sin}$ 2 ${\rm ß}$ -J $\rho$0 2 一 r(r+ 1)儿1a $\tilde{sin}$ 2 Ll 卢i r 主L l\textbackslash\{\} A\textasciitilde{} 2 \textasciitilde{};\textasciitilde{} 2 [J一芝工11tT P OI - Lz 十 (r - l) 地 i sinWMITU-l … p y干lJ (7.53) (7.54) (7.55) (7.56) 利用式 (7.51b) ，可导出类似于正激波的普朗特关系式如下(只需将总熔 h o 换 T r2 成 h. \textasciitilde{} = h o 一 亏，然后用前面相同的推导方法可得.请读者自己完成证明) U,..U... = 旦Z二，!.2.h. \textasciitilde{} = 主主二卫队。一旦 ${\rm i}$ = c2${\rm •}$ - r二$\tilde{U}$\textasciitilde{}' ${\rm ‘}$ r+l " V r+l 飞 " 2/ - . r+l 所以斜激波前后有 U ,..U... = c2${\rm •}$ - r二$\tilde{U}$$\tilde{cos}$2B 川 $\omega$ r+l-'-- - "" (7.57) 以上各式说明，斜激波前后气流 参 数的变化不仅与 r ， Mal 有关，还与激波角度卢有 关.激波倾斜角度卢与 气 流偏转角 S 的关系式可由下面几何关系导出. 因 叫 = 号， 消 去 U ， .Jt-代入式 (7.53) 可得 tan(${\rm ß}$- 8> =号 一主一+陆$\tilde{sin}$2 ${\rm ß}$ tan(${\rm ß}$ - 8) \_ U2n \_ Pl \_ r-l tan卢 U 1n P2 r土1 地$\tilde{sin}$ 2 ${\rm ß}$ Y 一 l 应用 三 角函数公式和代数运算，可以导出以下公式: n 儿1a $\tilde{sin}$ 2 $\beta$ - 1 tan8 = 2cot${\rm ß}$ ${\rm •}$ -:-:--:i ,... 1\textbackslash\{\}企l \textasciitilde{} ( r + cos2${\rm ß}$) + 2 这就是激波倾斜角与 气 流转角间的关系，并示于图 7. 15. 利用图 7.15 和式 (7.58> 进行分析，可得到斜激波的 一 些性质. (7.58)

48 。 44. 40. 36. 32. 咱 28 . 24. 20 。 16. 12" 8. 4.

\subsection*{例 0.20  . 30. 40. 50. 60. 70. 80. 90.}

卢 图 7.15 斜激波后的气流转折角和激波倾斜角间的关系 ${\rm 1}$当 Ma$\tilde{sin}$ 2 ${\rm ß}$-l =0 时， ${\rm ß}$=arcsin( I/Mal) = $\mu$ ，恰好是马赫角，这时 tan8=0 ， 这说明气流转角等于零时，斜激波退化为微弱扰动波，并称之为马赫波. ${\rm 2}$当 cot${\rm ß}$=O 时，卢=$\pi$/2 ，这就是正激波的情况，这时气流转折角也等于零. ${\rm 3}$式 (7.58) 中阳的是双值函数.给定 Mal' 卢和 $\delta$ 之间的关系可能出现 三 种情 况:(j)只有 一 个卢解，它对应于给定 Mal 时的最大气流偏转角 8m. .; (ii) 气流偏转角 大于 8m a .H才，卢没有解; (iii) 气流偏转角小于 8mB ，时，有两个卢解，其中 一 个解Ma 2 <1 ， 称 为强解，强解只在紧靠 8n旧的 一 小部分区间;大部分情况，斜激波后 Ma2 > 1. 称为弱解. ${\rm 4}$对式 (7.58) 求导数，可以求出对应于 Mal 的最大气流偏转角 8ma ${\rm ••}$ 超过这个角 度，满足斜激波关系的解卢不存在.而且 Mal 越大，最大转折角 8ma ${\rm •}$ 越大(参见 图 7. 15). 极限转折角对应的物理现象是斜激波由附体到脱体的转变.例如，在固定 Ma] 的超声速气流中放置一个可变顶角的尖棋，当模顶角小于 8ma x 时，在模顶处产生

附体斜激波(图 7. 16(a)); 当模顶角逐渐增加时，激波倾斜角随之增加，一旦模顶角 增加到 8ma ${\rm •}$ 时，再有微小的模顶角增加，斜激波会突然离模顶，在模顶前形成一道脱 体曲激波，如图 7.16(b) 所示.当超声速气流绕钝头物体流动时，也会在钝头物体前 产生脱体激波.总之，当斜激波后的气流转折角大于相应马赫数的最大气流转折角 后，激波不能稳定地附着在壁面，而要向前推移，使间断面脱离物面形成脱体激波.

\[MUl>1\]

(a) (b) (c) 图7.16 脱体激波的情况 ${\rm 5}$ Ma. 趋于无穷大时最大气流转角是附体斜激波的最大转角，对于 y= 1. 4 的 气体，其最大气流转角 8mo ${\rm •}$ = 45. 4 0 ${\rm •}$ 就是说，大于 8m". = 45. 旷的附体斜激波是不存 在的. 超声速气流通过斜激波时的气流参数的计算基本上和气流通过正激波时相同， 但是斜激波比正激波多一个参数 z 激波倾斜角卢或气流转折角 8. 这两个参数之间的 关系由式 (7.58) 给出.计算斜激波过程的气流参数时，首先利用式 (7. 58) 确定激波倾 斜角 ${\rm ß}$. 然后算出 Ma...= Ma l si 咐，即垂直于斜激波间断面的马赫数，余下的计算就同 正激波过程的计算一样.

\subsection*{例 7.7  设有马赫数 Ma=3 的超声速气流(空气) .流过尖棋，它的半顶角等于}

10 0 ，来流的压强 $\rho$. = 100kN/m 2 ，温度 T 1 =300 K.计算模面上的气流速度、压强和 温度. 解 首先应用式 (7.58) 或图 7.15 ，已知 Ma.=3 和气流转折角 8=10 。算出激波 倾斜角卢=27.4 0 计算 Mal ，， =Mal 町咐=1. 38 ，由正激波函数表(附表 5) 得斜激波后的气流参数为 岛也 2 . = 0.748 , $\rho$ZI$\rho$1 = 2.06 , TzlT. = 1. 24 最后得模面上的气流速度、压强和温度为 如= 2.06$\rho$1 = 206kN / m2 ${\rm •}$ T2 = 1. 24 T1 = 372K Ma 2 = 品公 z ，， /sin( 卢 一 的= 2. 50 U2 = a 2 l\textbackslash\{\}也 2 = 2.5-/ 1. 4 X 287 X 372 = 966(m/s)

(3) 斜激波的相交与反射 ${\rm 1}$固壁反射 第 7 章气体动力学基础 如果超声速气流中的尖棋可延伸至无穷远，则附体斜激波也将延伸至无穷远.实 际上尖模总是有限体，并且在流场中还常有其他固体壁面.当斜激波遇到其他壁面 时，在斜激波和壁面交点处，波后的气流方向必须和壁面平行(理想流体的边界条 件) .于是入射的激波在壁面上发生反射.斜激波和图壁相交可能发生图7. 17 所示的 三种情况.第一种情况，如图7.17(a) 所示，当斜激波和固壁相交时，交点后同壁的方向 和斜激波后的气流方向恰好平行，这时气流平滑流过同壁，这种情况下激波没有反射. 第二种情况(见图 7.17(b)): 假定固壁和斜激波前的气流平行.当斜激波和平壁相交 时，气流不能转折，于是在交点处就会产生第二道斜激波，迫使斜激波后的气流转到平 壁方向.假定来流的斜激波是弱激波，因此激波后气流仍是超声速.但是 Ma 2< Mal. 如 果第二次转折的转折角 ${\rm O}$2 小于J(;f应于 Ma 2 的最大转折角.这时产生的第 二 道斜激波是 由交点发出的直线.这种现象称为斜激波的规则反射.第一道斜激波称为人射波.第 二 道波称为反射波.第 王 种情况是不规则反射，也称马赫反射，如图7. 17(c) 所示，这种情 况是:入射斜激波遇图壁后壁面附近出现有 一 段垂直壁面的正激波，它和入射斜激波 相交，并在交点处反射一道斜激波，垂直于壁面的 一 段正激波称为马赫杆.当入射波后 气流和同壁的夹角大于对应 Ma ， 的最大转折角时.就会发生马赫反射现象. , ' , ${\rm ‘ð}$2

\[Mul >1 Mul>1\]

(a) (b) (c) 图7.17 斜激波的反射 (al 龙反射; (bl .lE 规反射; (c) '11 赫反射 通过以上分析可见，计算斜激波在图壁上的反射过程，实际上是已知气流参数 (人射斜激波后的气流参数)和l 气流转折角计算新的斜激波倾斜角和波后气流参数.

\subsection*{例 7.8  在平直槽道巾的超声速气流 Ma=3. 它的下壁面下游有 一 个半顶角 δ =}

10 。的尖模，它产生一道斜激波，计算该斜激波在上壁面反射后的马赫数和气流参数. 已知:来流的压强 $\rho$1= 100k :'IJ/ m2 .温度 T 1 =300 K. 解 本题的反射情况如图 7.17(b) 所示.斜激波后的气流参数已在例 7.7 中 求出 $\rho$2 = 2.06 户 1 = 206kN/m 2 ${\rm •}$ T2 = 1. 24 T1 = 372K. 1\textbackslash\{\}$\iota$ J'y生 2 . 一= 2.50 sin(${\rm ß}$ - ${\rm o}$) 已知第一道入射披后的转折角是 $\delta$= 10.. 反射波后气流平行于壁面，因此反射波的气

\section*{7.5  定常超声速气流绕凸角流动\{普朗特·迈耶流动)}\addcontentsline{toc}{section}{7.5 定常超声速气流绕凸角流动\{普朗特·迈耶流动)}

流转折角也是 8= 10 0 ${\rm •}$ 由斜激波后 Ma=2.5 和转折角 8= 10 0 ${\rm •}$ 通过式 (7.58) 或图 7. 15 可以求出反射披的激波倾斜角 ${\rm ß}$Ut = 30 0 ${\rm •}$ 反射披前的法向马赫数 M$\alpha$1 ，. = 2. 5sin30 0 = 1. 25. 由正激波关系式(或附表白，可得反射后的压强比、温度比 z $\rho$反射 / P2 = 1. 656. T反射 /T2 = 1. 159 最后，可以求出反射披后的压强和温度为 P ，i. 时= 1. 656$\rho$ 2 = 341kN / m2${\rm •}$ T眨射= 1. 159T 2 = 431K ${\rm 2}$强度相同的两斜激波的相交与反射(见图7. 18) 在实际气流中还可能发生两道斜激波相交的情况，如图7. 18 所示.这时相交的 两斜激波必然反射另外两条斜激波.因为每一道人射斜激波后气流都发生转折.但是 转折角不同，两股气流汇合后.只能有一个共同的方向，从而在两入射波的交点处反 射两道新的斜激波.如果来流的两道斜激波的强度相等(图7. 18(a)). 那么汇合后的 气流方向应和来流方向相同.利用斜激波关系式也可以定量分析这 一 过程，定解条件 是反射激波后气流速度大小、方向相同，压强相等(户 I$\rho$5 ). (a) (b) 图7.18 (a) 两强度相悖的斜激波相交; (b) 两强度不相等的斜激波相交 两道强度不相等的入射斜激波相交(图 7. 18<b)). 这时反射激波后的气流速度 方向必须相同，但是和来流方向不再平行.反射后${\rm 4}$、${\rm 5}$两 l叉的压强应当相等，${\rm 4}$、${\rm 5}$ 两区的气流速度方向相同但是大小不等，因而在${\rm 4}$、${\rm 5}$两区间存在切向的间断(又称 滑移线) .如图7.18<b) 所示. 在关激波和固壁、激波和激波之间的相互作用在气体力学专门课程将有详细讨 论，本书从略.

1.流动现象 前面讨论过超声速气流绕内折角的流动，即气体压缩性偏转产生斜激波.超声速

气流绕过凸角流动会出现什么现象呢?本节将专门讨论这一情况. 二工 Ma>1 外折角流动 连续膨胀 图7.19 超声速流的绕凸角流动 在斜激波理论中已经指出，激波偏转角很小的斜激波可近似为马赫波.假设超声 速气流绕微小外折角的平面流动用弱斜激波公式来估算(这时转角为负值) ，结果得 到激被后的压强小于激波前的压强，即负转角的斜激波是膨胀过程.前面已经论述 过:绝热膨胀过程不可能发生间断.闵此超声速气流绕凸角流动 一 定是连续等崎膨 胀.无论是绕外折角或绕连续凸面的流动图 7.19. 气流都是连续地发生转折和膨胀. 普朗特和迈耶最早研究这种超声速流动现象并得到解析解，因此这种流动又称为普 朗特-迈耶 (Prantdl- Meyer) 流动. 2. 普朗特-迈耶流动关系式 下面分析连续转折的超声速气流运动.并取如图 7.20 所示的控制体.均匀气流 在某一直线上开始膨胀，并发生转折，这种转折产生的扰动以马赫波向下游传播.也 就是说超声速气流绕凸角的平面流动是通过一系列连续转折的马赫披完成等情膨 胀，因此连续转折的马赫波又称膨胀波.根据上述分析首先导出超声速气流通过 一 道 马赫波微弱膨胀的普朗特-迈耶流动关系式. 皮厂 k 胀一阳、膨

\[-F\]

「飞 4 节''、，，， 向$\upsilon$ 4

\[: -14--E\]

图7.20 推导普朗特-迈耶流动用阁 取如图 7.20 所示的控制体，设气流在马赫披上转折前后的夹角为 d$\alpha$( 以 U+dU 逆时针转向 U 为正值，参见图 7.20). 根据动量方程可知膨胀披前后平行于马赫线的 速度切向分量相等，故有 Ucosp = (U + dU)cos( $\mu$+ da ) \textasciitilde{} (U + dU)[co年 -sm$\mu$d$\alpha$]

整理后得 $\tilde{Uco}$年 +dUco年一 Ud$\alpha$sm$\mu$ da = $\tilde{l}$! cotu = $\tilde{l}$! . /匕坐z\_p\_$\alpha$= 万 COI$\mu$= 万吨J sinz$\mu$ 已知马赫角和马赫数间有关系式 : sin.$\mu$ =l/Ma. 因此(图 7.20 中也是顺时针方向转 向，因此应取负值) 厅，....一-， dU da =- v'\textasciicircum{}也 Z 一 1 万 (7.59) 式 (7.59) 给出了等铺膨胀过程的微元转折角和马赫数以及速度相对增长率间的关 系.进一步，由能量方程消去 dU/U. 就可得到转折角和马赫数间的函数关系.等情过 程的能量方程可表示为 To \_ c\textasciitilde{} \_ 1 I Y 一 l七=斗 =1+ 一$\tau$一地 2 1 C \textasciitilde{} dc \_ c 1 y-1 .., z \textbackslash\{\} 1, I Y 一 1 ， "-2\textbackslash\{\} - 1将上式对 Ma 求导数.得$\cdot $一一一一一(一一比1( 1+ 一$\tilde{Ma}$21 ${\rm •}$ 由该式进一步. dMa Ma 飞 2 \_.\_- J 飞 2 \_.\_- J 可导出 dU/U 如下: dU d(cMa) d品也 r" dc 1 c \textbackslash\{\} -1 l 一--一一一一=一一 11 +一:: 1.-7- 1 1 U cMa Ma L- , d 此也飞 Ma J J =警 [1 一(早地 2) (1 +早地 2fJ =带 (1 +早地 2fl 代人式 (7.59) 得 积分上式得 da 一-\_j\_商$\tau$$\tau$T dMa n 一- u 叩 -1" ? Ma 1+\textasciitilde{}一岛也 Z = ${\rm •}$ /吐l arcurljZ二\textasciitilde{}(Ma 2 - 1) - arctan ./Ma 2 可 +c 匀 y-l 勾 $\gamma$+1 令 Ma=l. $\alpha$=0 ，可确定常数 c=o. 由此得到的气流角称为普朗特-迈耶角，用 $\nu$ 表示 z $\nu$( 协) = . /过i arctant1( 地 2 - 1) 一川an /Ma 2 可 叫$\tilde{y}$-l 勾 y 十 l (7. 60) 式中 $\nu$ 的物理意义是:气流从 Ma=l 膨胀到 Ma 时的气流转折角.该关系式已做成 图表(见附表 6). 由式 (7.60) 可得到超声速等蛐膨胀转折的一些结论. ${\rm 1}$当 Ma $\rightarrow$$\infty$， $\nu$ 二 $\nu$nlll. X -$\pi$( / (y+ 1) / (y-1) -1 ) /2. 这是超声速等情膨胀的最 大转折角.对于常见双原子气体 .y= 1. 4.$\nu$m.. = 130.5 0 ${\rm •}$ 理想气体等情膨胀到 Ma $\rightarrow$$\infty$ 时，压强和绝对温度均为零，所以最大转折角可解释为 Ma=l 的超声速气流绕一半

无限长薄平板流动，薄板的下部压强等于零(图 7.2 1) ，这时气流绕过平板边缘，它的 最大转角为 130.5 0 ，大于 130.5 。 处将出现压强等于零的真空区. ${\rm 2}$式 (7. 60) 中 $\nu$ (Ma) 是单值函数 ， Ma 从 1 增加到无限大， $\nu$ 值单调地从 0 增大 到 130. 5 0 (y= 1. 4). 若气流绕凸壁面连续膨胀，已知凸面起始切线和终点切线的夹 角，即气流的转角，可利用以 Ma) 函数式(或附表 6) 直接通过已知的 Ma] 和转角求出 Ma 2 或已知 Ma] 、 Ma2 求出气流转角.具体算法见以下例题.

\[Ma=1\]

真空区 Mu] 图7.21 超声速 气 流的最大等娟转角 图7.22 超声速 气 流绕 二 维凸形物流动示意图

\subsection*{例 7.9  已知 Ma] =2.0 的均匀超声速气流绕过 α = 7.9}

0 的 二 维凸形物面，如 图 7.22 所示，求绕过物形后的马赫数 Ma 2 . 解 由 Ma] = 2.0 ，查附表 6.$\nu$ (Ma) 函数表，得 $\nu$] (2. 0) = 26. 38 0 ，已知转角 $\nu$2 一 $\nu$]$\alpha$ = 7.9 0 ，故问 = $\alpha$ + $\nu$] = 34. 28 0 ，由 $\nu$2 值查附表 6 ，得 Ma z = 2.3. 就 是 说 Ma] =2.0 的超声速气流.绕凸角流动转过 7.9 0 后加速至IJ Ma z= 2. 3.

\section*{7.6  完全气体在变截面绝热管内的准一维定常流动}\addcontentsline{toc}{section}{7.6 完全气体在变截面绝热管内的准一维定常流动}

本章最后，综合应用前面讲述内容，分析和计算 一 个 r 程实用问题: 气 体在变截 面管道中的绝热流动.这是 一 种高速喷管流动，常用于火箭推进器、涡轮机的静叶片 和动叶片通道，以及超声速风洞等.作为 一 种 t 程设计方法，需要把实际问题简化成 易于分析计算的模型，同时又保留流动的主要性质.对于气体在变截面管道中的绝热 定常流动，仁程中常采用如下简化模型. 1.准一维定常绝热假定 气体在变截面通道中流动时(如图 7.23 所示) ${\rm •}$ 假定流动参数在同 一 截面上是均匀的，只在流动方 向发生变化，这种近似模型称为准 一 维假定. 一 维变 截面流动中，流动的边界条件是通道截面积 A = A( 抖，其他流动参 量 都是 r 的函数，如 : u=u( 川， \textbackslash\{\}y$\gamma$/ \textasciitilde{}飞亏、 r 图7.23 气 体的 一 维流动

p=p( .r) ， $\rho$=$\rho$ ( .r) 等.准一维流动近似忽略了流动的横向分量，如果流动的主流速度 远远大于横向速度，准一维假定是很好的近似.从几何边界条件来说:如果通道截面 的变化率很小，就能满足一维近似的要求.具体来说，准一维近似要求 1二兰主<<1，其. A o d.r 中 L 是通道的长度， A 。是通道进口截面积. 如果外界没有热量输入，气体流动过程也没有化学反应、蒸发等内部生成热，气 体的粘度又很小，气体流动可以认为是理想绝热的.下面给出准一维定常绝热流动的 分析和计算方法. 2. 准一维定常绝热流动基本方程 根据准 一 维定常流动假设，应用理想流体运动的基本方程，可导出以下准一维定 常绝热流动方程. (1)连续方程 (川) = 0 (7.61) 上式就是定常流动沿流管的质量守恒方程:通过任意截面的质量流量相等，即

\[puA = Q (7. 62)\]

(2) 运动方程 ，主 1 、。一吨。1-p 一一 句

\[-hu (7. 63)\]

(3) 能量方程 u  (h 十 二) = o. 或  (h+ 专) = 0 (7.64a) 上式就是理想流体绝热定常流动沿流管的能量守恒方程: h+ 亏= ho (7.64b) 对于完全气体 $\tilde{P}$\_+ 旦二 = h" (7.64c) Y 一 1 p ' 2 .." 若流动是连续的，也就是没有激波间断时，绝热流动应是等峭的，因此还有 生 y 户 一一 h凹-y 户 3. 气体准一维定常等煽流动的主要性质 (1)气体准 一 维定常等情流动中，流速与截面积变化之间有如下的关系式: ? "du dA (Ma ' 一 1 ) 一 =万 (7. 65)

证明 将连续方程 (7.61) 展开，可得 dp I du I dA 一+一 +$\tau$\textasciitilde{}=Op U /"i 将 c 2 =(d$\rho$/dp) ，代入动量方程 (7.63) ，可导出 pudu =- dp =一 (dp/dp).dp =- c2 dp 或 dp =\_ u\textasciitilde{} 但 =-Ma 2 生 p c u u 将上式代入微分型的连续方程后，可得式 (7.65) du dA 仙也 2 -1)一

\[U /"i\]

证毕. 由式 (7.6 日，可得到气体准一维定常绝热连续(即等摘)流动的以下性质. ${\rm 1}$气体的亚声速流动中，即 Ma<1 时，若 dA>O ，则 du<O; 反之，若 dA<O ，则 du>O. 就是说气体亚音速一维等摘管流中，截面积增加，流速减小 z 截面积减小，流 速增加.这一性质和不可压缩流动类似. ${\rm 2}$气体超声速流动中，即 Ma>1 时，若dA >O ，则 du>O; 若 dA<O ，则 du<O. 这就意味着 z 气体超声速一维等摘流动中，截面积增加，气流加速 z 截面积减小，气流 减速.这一性质与不可压缩流动完全不同，这是气体超声速流动的固有特性. ${\rm 3}$ Ma=1 时，根据式 (7.65) 可导出: dA=O. 就是说气流中出现声速 (Ma=1 )的 截面是截面积变化的极值点.在等湍流情况下还可进一步证明它是截面积的极小值， 即管道内流动若达到声速，一定在最小截面处. 根据以上分析，可定性地了解气体在变截面通道内可能的等情流动. (2) 收缩通道流动 亚声速定常等情流在收缩通道中将加速，但始终保持亚声速\$超声速定常等辅流 在收缩通道中将减速，但始终保持超声速.如图 7.24 所示.

\['a M >i a >, ,'\]

叫 M

\begin{gather*}
,-U M <1\\
1\alpha \\
l <U M
\end{gather*}

图 7.24 气体在收缩绝热通道中流动示意阁 (3) 扩张通道流动 亚声速定常等惰流在扩张通道中将减速，并保持亚声速 z 超声速定常等情流在扩 张通道中将加速，且始终保持超声速.如图 7.25 所示

\[MlI, <1 Ma2>MlI, >I\]

图7.25 气体在扩张绝热通道中流动示意阁 (4) 收缩扩张管流 收缩扩张通道中气体等情流动情况较简单收缩或扩张通道中流动复杂. ${\rm 1}$亚声速定常气流进入绝热收缩扩张通道时，在收缩通道中亚声速等'脑气流必 加速，而后可能发生以下 三 种情况: (i) 先在收缩通道中加速，当气流到达最小截面 时.流动仍是亚声速，这时问:进入扩张通道后气流将减速，全部通道巾是亚声速流，如 图 7.26(a) 所示; (ii) 亚声速气流开始在收缩通道中加速，在最小截面处达到声速，但 是通道出口压强较高，气流进入扩张通道后又将减速，而成为亚声速，图 7.26(b) 表 示这种情况; (ii i) 亚声速气流在最小截面处达到声速，而通道出口压强较低，那么气 流进入扩张通道后将继续加速而达到超声速，如图 7.26(c) 所示.上述分析表明亚声 速气流在收缩扩张通道中的流动情况，取决于通道出口压强和通道最小截面上的马 赫数.

\[dA=O.MlI <1 dA=O.Ma=1 dA=O.MlI =1\]

MlI, <1 1 、中=卢卢-1 1 气毡，占户户 1 1 飞也卢户户 l -一\_\_ I 一- ， 一一- 一-- -一一-… --<… 一+-一一一… -$\rightarrow$一-一… $\nu$'户 ' 才亏$\tilde{Ma}$,<1 V" '/， ，/巧节，. MCl2< 1 V$\gamma$ / 巧乍$\tilde{Ma}$\textasciitilde{} > 1 (a) ) LU ( (c) 图7.26 入口亚f.1速气流在绝热收缩扩张通道中的流动 ${\rm 2}$考察超声速定常气流进入绝热收缩扩张通道时的情况，超声速气流在收缩通 道 'P 必减速，而后也可能发生 三 种情况 : (i) 先在收缩通道中减速，如果气流到达最 小截面时，流动仍是超声速的，进入扩张通道后气流将变为加速.如果出口压强较低， 则全部通道中可以是超声速流，图 7.27(a) 表示这种 t 况; (ii) 如果超声速气流在最 小截面处达到声速，但是通道出口压强较高，气流进入扩张通道后继续减速.而成为 亚声速，图 7.27(b) 表示这种情况; (iii) 超声速气流在收缩通道中减速，如果气流在 最小截面处达到声速，而通道出口压强较低，那么气流进入扩张通道后将加速而达到 超声速，如图 7.27(c) 所示. 如果在扩张通道中出现超声速流动，是否能够继续维持超声速流，要视通道出口 压强而定，后面将有例子说明，扩张通道中的超声速流动可能在通道内或通道出口产 生激波.

\[dA=O.Ma>1\]

Ma ,>1 I\textasciitilde{}卢---1 d4 =O. 儿1<1 = 1 dA = O. 儿1<1 = 1 $\tilde{I}$ \textasciitilde{} 1 ---1… …--唱… .... …- --+…--<…-+-… 1,;""/ ""，'7女":'7泸 Ma2> 1 .,....". " // \textasciitilde{}月" Ma、 <1 V;-" -" -'777为4 品1，$\alpha$ ， > 1 (a) (b) (c) l羽 7.27 人口超声速气流在绝热收缩扩张通道中的流动 4. 流量公式与密流性质 (1)流量公式 根据连续性条件，变截面管道内气体定常等摘流动的质量流量在各截面处相等， 将气体等情流动关系式代入流量公式.可得 ,n = puA = A 丘ptJ MJ Lco= A而(J 问 (1 + 与土地 2 ) 2c y- " ( 7. 66) p o Co (2) 流量密度和性质 定义 7.11 单位面积上通过的质量流量称作流量密度.简称"密流"，用 q 表示. 即 q=pu. 根据密流定义，从前面流量公式 (7. 66) 可得流量密度与马赫数的关系式: ....L土L 4 儿也) = pu = 布:;五汕 (1 +号..!Ma 2) '(1 11 气体定常等恼流动中密流有以下性质. ${\rm 1}$在相同的滞止状态下 .Mu=l n才密流最大. 由 dq(Mu)/dMa=O. 可得 Ma = l 为函数 q(Ma) 的极值点，而且可证明 : Ma = l 时 d 2

\[q / dMa\]

2 < O. 网此 Ma=1 时密流为极大值.图 7.28 展示密流和马赫数的关系. 从图 7.28 可以看到，在定常等恼流中.当 Ma < l(即亚声速)时，密流随马赫数的增加而增大，因此 在给定质量流量的管流中，马赫数增加(加速)截面 积必须减小;而在超声速流动情况下，密流随马赫数 的增加而减小，因此当马赫数增大时(加速)截面帜 必须增大. ${\rm 2}$临界截面 A ${\rm ‘}$与墙塞流量. 定义 7. 12 在变截面等蛐定常流动中 Ma = l 的截面称为"临界截面"，用 A' 表示. 图7.28 气体等情流中密流和l 马赫数间的关系 从前面的分析不难得出结论 z 在气体变截面定常等摘流中，临界截面一定是最 小截面，工程中常称最小截面为喉部.如果等恼定常气流在收缩扩张通道的最小截面 处达到 Ma=l ，则根据密流极值的性质，此时通过收缩扩张通道的流量是给定滞止 参数下的最大流量.称之为堵塞流量.最大流量可由下式算出 z

p' u' A ' = A ${\rm •}$ p' c' = A ' .j r${\rm Þ}$ $\rho$ 。 ($\tilde{ll}$ )琦J日 (7 ， 67)V TP opo \textbackslash\{\} r + 1 \} 堵塞流量是给定滞止参数下，变截面管流中气体等摘流动可能达到的最大流量.就是 说，当滞止参数给定后，降低管道出口压强，通过管道的气体流量不断增加;当流量达 到墙塞流最后，再降低出口压强，通过管道的流量不会再增加.根据变截面气体的等 情流动原理.这时在管道的最小截面上，气体流速达到当地声速，而这 一 速度是气流 在最小截面上的最大速度，无论怎样减小出口压强，都不会使最小截面上的速度 增大. ${\rm 3}$定常等铺流中截面积与马赫数的关系. 由连续方程: puA = p' u' A' ， 并将流动参数的等蛐关系式代人，可导出以下关 系式: A D ' ll ' 1/2 ， y 一， 、 古」T丁 一一 =$\iota$二.一:.-(一一一 + $\tilde{Ma}$ " l A' pu 1\textbackslash\{\}也飞 r + 1 . r + 1--- I 上式用曲线表示于图 7 ， 29 ， 由图 7 ， 29 可看出， 一 个 A / A' 对应于两个 Ma 值，其中 一 个 Ma<l ， 另一个 Ma > l. 具体流动过程 中对应哪个值，由流动进出口边界条件确定.正如前 面收缩扩张流动中讨论过，当喉部截面为临界截面 时，扩张段可能出现越声速流动.也可能是亚声速流 动.取决于扩张段出口的环境 ffi3虽. (7 , 68) 45 40 35 nu$\epsilon$Jnu$\epsilon$J 句37-7--l 飞\textbackslash\{\}T 10 (3) 收缩喷管的 t 程计算 利用流蛙公式( 7, 66) 和等'脑流的参数关系 图7.29 式 (7 ， 25)- 式 (7 ， 33) 可以计算气体准一维定常等饷 吨赫数间的关系 流动.在气体的喷管流动中，喷管出口截面压强用 $\rho$， 表示.喷管 H-: U 环境 ffi 强用 $\rho$" 表示， $\rho$， 可以等于 $\rho$ /， 也可以不等于队，详见下文分 析.通常有两类计算问题:设计问题(也叫反问题)和运行问题(亦称正问题) ，分别介 绍如下. ${\rm 1}$设计问题 z 给定流量 m 、滞止参数 $\rho$ " 、 p" 和出口马赫数 Ma(< l)， 求出口截面 有1 A . 和怀.强 $\rho$ . ' 求解步骤如下. 第 一 步:利用等情关系主 (Ma) ${\rm •}$ 由给定马赫数计算出口压强户. ; 第 二 步:用流量公式 (7 ， 66) 计- 算出口截面积 A . ， 在设计问题中坯可能给定流 量 m 、出口马赫数、出口气体参数，要求计算滞止参 数和出口截面帜.设计步骤如下.

第一步:利用等情关系主 (Ma) ， 由给定马赫数计算滞止压强; 户。 第二步 z 用流量公式 (7.66) 计算出口截面积 A.. ${\rm 2}$运行问题:已设计好的喷管，在实际运行时出口压强|坷环境状态而改变，这 时通过喷管的流动和设计状态不同，需要重新计算.运行问题的提法是: 给定出口截面积人、滞止参数扣响和出口环境压强扣，要求计算通过喷管的流 量 m 和出口马赫数 Ma. 求解步骤如下. 第一步:计算户1 / 扣$\cdot $并用等'脑临界公式判断出口截面是否为临界状态? cl. 若出口截面 $\rho$$\rho$ p' ， 则出口马赫数必小于 1 .根据等情流公式手 (Ma) 计算出 p 。 口 Ma.; b. 若出口截面如 $\zeta$ p' ， 则出口截面必为临界截面.因为收缩绝热通道中等恼气 流最大马赫数等于1，并且出现在通道的最小截面处.因此当如 \textasciitilde{}${\rm þ}$' 时，出口截面上 的压强 $\rho$，等于等摘流的临界压强户 ，. =p* ， 而口截面上的马赫数 Ma.= l. 下标 r 表示 出口截面的状态. 第二步 z a. 若出口 $\rho$，， >$\rho$$\cdot $. 用式 (7.66) 计算流量 m; b. 若出口为临界状态，用式 (7.6 7)计算流量 m; C. 岛主二 $\rho$$\cdot $ 时，气流将在管外膨胀. 在上面分析中已指出，在运行问题中，当环境压强小于临界压强时(户b< $\rho$ ${\rm ‘}$) ，气 流在管内只能膨胀到临界压强，因此出口截面上的压强大于出 U 外的环境压强.气体 流出喷口后，还要继续膨胀，这种情况称为膨胀不足.如果喷管流动是平面等崎流动， 气体在出口外膨胀可以通过普朗特 - 迈耶流的绕角流动来分析计算，由等情流公式 主主 (Ma) 计算气流完全膨胀的马赫数，然后通过普朗特-迈耶函数关系式 $\nu$ (Ma) 计算 p 。 气流出口后的转折角.此时，收缩喷管出口外气流继续膨胀，并发生转折，因此流动不 再是一维的，如图 7.30 所示.平面喷管出口外的流动比较复杂，其基本过程如下:气 流经普朗特-迈耶流动膨胀后，气流边界呈扩张型;但是出口上下边缘发生等孀膨胀 产生的超声速气流的方向是不同的，于是 在出口上下缘发出膨胀波，两膨胀波在中 心相交穿过，并保持为膨胀波，其自由面向

\[Mu<1\]

外扩张，与自由边界面相遇后反射压缩披， 自由面收缩，两压缩波相交反射压缩波，再 图7.30 收缩喷管的口外膨胀 遇自由边界面后反射膨胀波，气流在出口 细实线 2 气流边界 5 粗实线 g 激波 s 虚线 z 膨胀披 后经历反复膨胀和压缩直到气流演化为亚

声速射流.以上过程部分地示意于图 7.30. 详细的计算方法可以参考气体力学书籍. 5. 拉瓦尔( LavaD 喷管内定常绝热流动 前面已经论述过，气体只有通过收缩扩张通道才能在通道出口处达到超声速，这 一节分析收缩扩张喷管(也叫拉瓦尔喷管)的设计(反问题)和运行问题(正问题). (1)设计问题:给定流量 m ， 滞止参数 $\rho$。 、户和出口马赫数 Ma(> l)， 求最小截 面积 A ' 、 A 俨和出口压强户 ， .求解步骤如下. 第 一 步:利用等恼关系主 (Ma) ${\rm •}$ 根据给定的出口马赫数 Ma(x) 计算压强 $\rho$， ; 第二步:用流量公式 (7. 66) 计算截面积A\_; 第 三 步 z 准确计算最小截面积(常称喉部截面积) A" = m/(p. u' );或由 式 (7.68) ，利用出口马赫数和出口截面积计算喉部面积. (2) 运行问题.按等摘超声速定常流动条件设计的拉瓦尔喷管，当环境压强发生 变化时，喷管内的流动会发生什么变化，这是运行问题.需要求解的问题是 z 给定:气体的精止参数 $\rho$。 、 T" 、喷管的临界截面积 A" 、出口截面积 A ，和出口环 境庇强(背压)如.求通过喷管的流量和出口马赫数. ${\rm 1}$几个特征压强. 收缩扩张喷管中的超声速定常绝热流动比较复杂，在讨论变工况的流动之前，首 先分析几种典型流动工况喷管出口截面的压强，即特征压强(见图 7. 31). 出口截面的流动参数用下标 "r" 表示，环境压强用户，表示，喉部截面流动参数用 下标 "1" 表示. 第 一 种典型流动 T 况是:管内完全是等摘亚声速流动，喉部 Ma ， =l ， 收缩段和 扩张段都是亚声速流动，这时出口压强用 $\rho$${\rm ‘}$表示，如图 7.31 上的${\rm 2}$; 第二种典型流动 I 况是:管内完全等情流动，喉部 Ma ， =l ， 扩张段为超声速流 动时的出口压强用 $\rho$J 表示，如图 7.31 上的${\rm 6}$; 第 三 种典型流工况是:出口截面前流动完全等惰，刚好在出口截面上产生一正 激波，激波前压强等于丸$\cdot $激波后压强用 $\rho$ f 表示.如图 7.31 上的${\rm 4}$. 下面用以上三个特征出口压强将流动分为 I 、 H 、阳、凹四种 t 况. ${\rm 2}$流动分析. T 况 1 :扣， / $\rho$。〉 $\rho$$\iota$ / 扣，喷管内全部是亚声速流动，此时 Ma ， <l .Ma ， <l. 出口 压强 $\rho$，扣，通过喷管的流量由环境压强确定，也就是说，对于这种士况，通过喷管 的流量对背压的变化很敏感. t 况 11 :$\rho$${\rm ‘}$ / $\rho$。〉 户， / Po> 扣/扣，即出口截面压强大于岛而小于丸，这时在管内 出现正激波，如图 7.31 中状态${\rm 3}$.正激波后为亚声速流动，在扩张通道中亚声速流动 是增压减速的，所以气流到达出口时 ， Ma， <l 且 $\rho$， $\rho$，， >Pr. 这种 t 况中背压舟，值越

Ma ${\rm 7}$飞\textasciitilde{}飞\textasciitilde{}-可气 ${\rm ø}$;::z.，;.-\textasciitilde{}勘U 图 7.31 典型的收缩扩张喷管出口工况示意图 高，正激波位置越靠近上游.在 t 况 E 中，喉部始终是声速，因此流量为常数，不受背压 影响. 工况田 z$\rho$ dPO >Pb/PO > $\rho$j/ $\rho$。，整个喷管内流动为超声速流动，因如〈岛，出口 截面不发生正激波，而且出口截面上的压强 $\rho$.$\rho$j<Pb( 这种工况又称膨胀过度). 这种工况下，管内流动不受如变化的影响.由于 $\rho$， -$\rho$j< 冉，气流流出喷管后发生压 缩.我们知道超声速压缩会产生激波，现在已知激波前后的压强比等于血战.>1. 激 波前的马赫数等于 Ma ，. 应用激波理论由压强比和马赫数可以确定激波的特性，如 2 .激波倾斜角和激波后气流参数等. 工况凹: Ph/ $\rho$。 〈丸 /$\rho$。，喷管内的流动不受 Ph 变化的影响，出口截面上 $\rho$ . -$\rho$尸 管内全部是等情流动.由于 $\rho$ ， -$\rho$j> 冉，这种工况属于膨胀不足的情况，因此气流流 出喷管后将继续膨胀.对于平面喷管，喷管外的流动可以应用普朗特-迈耶流动和斜 激波理论分析并计算出口外的膨胀一压缩一膨胀一压缩......过程，其基本情况和收 缩喷管膨胀不足的情况相同. 以上讨论的各种流动T.况如图 7.32 所示. $\rho$h......\_ P ${\rm 1}$ 1> 一>\textasciitilde{} (工况 1 ) PO PO 取$\iota$$\mu$\textasciitilde{}气 Ma.< 1. 1> 去〉 Z Ma<1 、"'-""-户--

\[-\iota - Ma.<1 ~\]

\textasciitilde{}啊?听忏TJZb

${\rm 2}$生=主$\iota$ (如图 7.31 特征状况${\rm 2}$〉Po Po 胸怀$\iota$$\mu$$\mu$\textasciitilde{}$\iota$fMae〈 l $\rightarrow$-一$\tilde{a}$， =1 寸吃 =n =n$\tilde{Ph}$\_P c 扩俨 / / 厅开开?疗JIL61 …内局 $\rho$g$\rho$b....\_ P! ${\rm 3}$一〉一〉 $\iota$ (工况 10

\[p .  \rho  .  \rho  . \]

${\rm 4}$主L= 主主〈主二 〈如图 7.31 特征丁， 况${\rm 4}$〉Po Po -Po P!....\_ Pb 由${\rm 5}$ 1"'> 一 >rJ (工况皿)

\[Po \tilde{P} .  \rho  . \]

@生=主i... (如图 7.31 特征主况${\rm 6}$)Po Po 协〈半$\iota$$\mu$$\mu$A亏Z741

\[---l-- Ma. = 1 ---l-- -\]

$\nu$何市，..，..// 7))))i31h 去=去 $\rho$h \_\_\_\_Pj ${\rm 7}$一 <YJ (丁.况凹) Po $\tilde{Po}$ """乐创$\mu$$\mu$$\iota$/J JMa〉 ll;二' . .>1 -二.... 二 Ma.=1 \textasciitilde{}< \_ n W竹附节阴市7;二〉Pb 2〈去 图7.32 收缩扩张喷管中流动 t 况示意图

\subsection*{例 7.10  恒压容器中的空气由收缩扩张喷管流出，给定容器中压强 Po=5X}

1Q5Pa. To =288K. 喷管喉部面积为 10c 时，出口面积 20cm 2 ${\rm •}$ 试求以下出口环境压强 时，通过喷管的质量流量、出口截面的马赫数、压强、速度和出口后气流角度.

\begin{gather*}
(1 )\rho  1.=4.75XI0.Pa;\\
(2)\rho 1.=2. OX 10'Pa;\\
(3)\rho 1.=0.40XI0'Pa.
\end{gather*}

解由户。 ${\rm •}$ Tu 可计算向=户。 /(RT，，) =6. 05kg/m 3 ${\rm •}$ (1)先求特征压强丸，丸， $\rho$ l. 由 A 梅 /A.=10/20=0.5. 查等情表，得特征状态 Ma ， . . 丸 .M，饵， $\rho$j : 岛也2 ，- = 0.31. $\rho$ .1 $\rho$。= 0.9355. 队= 4.6775 X 10"Pa 此缸 j = 2.2. pJ $\rho$。= 0.09352. Pl = 0.4576 X 105 Pa 查正激波表，得特征状态 PI 去(地，= 2.2) = 5.48. p, = 2.508X 1Q5Pa (2) 判断给定背压下的流动范围，求解未知量. ${\rm 1}$扣 =4.75 X 105 Pa 时，如 =4. 75XI0 5 Pa>$\iota$ ，流动为全部亚声速流，所以出口 压强 $\rho$，如 =4.75XI0 5 Pa. 由 $\rho$b/ $\rho$。 =0.95. 查等捕流表得 品缸， = 0.27. T, /T" = 0.9856. Tr = 0.9856 X To = 283. 85K c. = I页T ， = 337. 7m/s. U. = c.Ma, = 9 1. 2m/s.

\[p. = \rho ,  /(RT.) = 5. 83kg/m\]

J

\[m = p,U,A. = 1. 063kg/s. 8. = 0\]

${\rm 2}$冉 =2. OX 10 5 Pa 时， $\rho$ ，>如〉岛，所以流动为管口产生斜激波的第皿种工况， 故由等情流动表，先算出出口截面的流动参数 z 儿缸. = 2.2. $\rho$. -岛= 0.4575 X 105 Pa T.lTo = O. 5081. T. = 146. 3K. c. = I页T. = 242. 46m/s

\[U. = c.Ma. = 533. 4m/s. rn = p.U.A. = 1. 188kg/s\]

出口后的气流参数继续计算如下 z 由出口背压如 /$\rho$j =2.0/0.4576=4.37. 查正 激波表可得 Ma'n =Ma , sin${\rm ß}$= 1. 97 ，因此卢=63.57 0 ${\rm •}$ 再由斜激波图 7. 15 查得气流转 角 8=26 0

\[( Pb=O. 4X 10\]

5 Pa 时，如〈丸，属 $\tau$. 况凹，为管外膨胀，而口截面参数由等楠表 计算 z 结果和${\rm 2}$相同- P.= $\rho$j = 0.4575 X 105Pa. 此缸， = 2.2. U.= 533.4m儿 ${\rm •}$ r${\rm n}$ = p.U,A. = 1. 188kg/s

\section*{习题}\addcontentsline{toc}{section}{习题}

管外膨胀为普朗特-迈耶流动，由等情表，仇 /$\rho$。 =0.08 时 Mab=2.3. 于是由普 朗特-迈耶函数:的 (Ma ， = 2.2) = 3 1. 73 0

\[.\nu 2 (Ma , = 2. 3) = 34. 28\]

0 ${\rm •}$ 出口气流转折角

\[8=34. 28\]

0

\[-3l. 72\]

0

\[=2. 55\]

0

\noindent\textbf{7.1.} 空气从储气罐由流，气罐内压力为1. 5 X 10; Pa. 温度为 288K. 已知出口为

亚声速流动，大气压力为 1 X 105 Pa ${\rm •}$ ;]<气流州口速度及温度.

\noindent\textbf{7.2.} 证明:理想完全气体定常等湍流动时有

(1)牛二主 =$\tau$L 「 \{l+$\zeta$1Maz 俨 -11

\[\div \rho  l..f ylVl a" L , t. I J\]

$\rho$。一 $\rho$1(2) Ma<<l 时， $\tau$一一 =1 十 -;-Ma 2 言p l..f $\tau$ 7.3 15km 高空飞行一枚导弹，测得表面温度 Tw 为 500K. 假定表面温度状态 阳驻点温度 M同，求导弹速度实际上旧是总温，应为 Tw=T+ 艺币.Pr 为普朗特数，对于空气 Pr=0.85. 那么导弹速度又为多少?

\noindent\textbf{7.4.} 空气在加热有摩擦的变截面管道内定常流动，已知某截面上 T=268K.

\[P=79. 2kN/m 2 .Ma=O. 6.A=6cm\]

2 ${\rm •}$ 求该截面的户IJ' To.p 饵. T' ${\rm •}$ c . ${\rm •}$ A ${\rm •}$ ${\rm •}$ U m.x ${\rm •}$ 7.S 空气等摘地通过收缩喷管流入压力为 124.5kN/ 旷的容器，喷管进口的空 气速度可忽略，进口压力和温度分别为 200kN/ 时和 293K. 出口面积为 78. 5cm 2 ${\rm •}$ 求 通过喷管的流量，又问该喷管的阻塞流量为多少?

\noindent\textbf{7.6.} 拉瓦尔喷管进口空气压力为 300kN / m

2 ${\rm •}$ 温度为 373K. 速度可忽略，管内 流动无激波.已知出口温度为 253K. 喉部面积为 10cm 2 ${\rm •}$ 求(1)出口马赫数; (2) 出口 面积; (3) 流量.

\noindent\textbf{7.7.} 直径为 0.2m 的水平管道输送压力为 825kN/m

2 、温度为 293K 的空气.拟 安装一文丘里流量计测量流量，流量范围不超过 llkg/s. 喉部压力要求大于 745kN/ m2 ${\rm •}$ 文丘里流量计的流量系数为平 =0.98( 实际需要面积=理论计算面积/流量系 数) .问实际需要的喉部直径 d ， 应为多少?

\noindent\textbf{7.8.} 空气通过渐缩喉管向大气放气如图 7.33 所示，大气压力为 100kN/ 时，进

口处 A ， = O. 558m 2 .$\rho$1= 150kN/m 2 ${\rm •}$ T, = 763K. 出口为亚声速 .A 2

\[=0.lm\]

2 ${\rm •}$ 求 (1)气流作用于喷管上的力; (2) 欲使喷管不动需加外力多少?

\begin{gather*}
PI---\\
T1 --- ~
\end{gather*}

图 7.33 题 7 . 8 示意阳 图7 . 34 题 7.9 示意阁

\noindent\textbf{7.9.} 海平面上空气通过图 7. 34 所示管道吸入真 空 箱，由纹影照片知 A， =

300cm\textasciitilde{} 处有一正激波，求真空箱内压力.已知 A T = 100cm2 .A ， =400cm\textasciitilde{}.

\noindent\textbf{7.10.} 总压管和总温度计放置在压力为 101kN / m

2 、温度为 288K 的静止大气 中，如前方发生爆炸，一股速度为 400m / 吕的热琅正对着它们袭来.问热浪袭来前后 总压管和总温度计读数各为多少?

\noindent\textbf{7.11.} 核爆炸波以波速 D= 15240m/s 在静止的空气Ij l 传播，静止空气压力为

\subsection*{例 01.33  kN/ m2 .温度为 294.4 K.试计算: (1)波后空气的绝对速度; (2) 波后\textasciitilde{}气相}

对于静止观察者的总压和总温.

\noindent\textbf{7.12.} 证明如图 7.35 所示两声波相碰后，原静止区中空气的压力增量 I:::. p 及 U

分别为

\[1 A L TT - 1:::. \rho 1- 1:::. \rho  2 jRT\]

I:::.p = 1:::. $\rho$1+ 1:::.扣. u 一一一--一!:..::. \textasciicircum{} j 飞子 p 'V "1 叫 十 p . T 十十十p俨+ 图 7.3汩5 题 7. 12 /J尽』民斗意阁 7.13 证明斜激波前后法向速度分量满足 U 1 .U 2 ${\rm •}$ = (" 飞一 $\zeta$1Y+1 P2 \_ 2"1 A${\rm ß}$'! "1- 1 7.14 证明斜激波前后压力有下式成立$\cdot $一 一 一'!--c- M 一一一: PI "1 +1"'1" "1+1' 7.15 顶角为 60 0 的尖模在 15km 高空作零攻角飞行，其飞行速度恰好使斜激波 附在尖模上，求: (1)飞行速度; (2) 激波角; (3) 尖模表面压力川的通过激波的总压 损失.

\noindent\textbf{7.16.} 用公式计算图 7.36 所示头部激波不被壁面反射的 $\alpha$ 值.已知来流的马赫

数为 3. 激波角为 30

Ma 国 =3 图 7.36 题 7.16 示意图 图7.37 题 7.17 示意图

\noindent\textbf{7.17.} 如图 7.37 所示，求 4 区与 5 区的气流马赫数、压强 $\rho$ 和方向角。.已知

Mal =4.0. $\rho$1 = 100kN/m 2 ${\rm •}$ 7.18 二维平板在 6km 高空以 Ma=2 飞行，攻角 10 0 ${\rm •}$ 弦长 c 为 1m. 求 z (1)上下表面的压强和速度; (2) 升力系数 C，. 和阻力系数 C D IC [l， cL 的计算公式是 : CL = 丁 FL 一 .CD = l 言p"， U$\tilde{x}$ - c

\[F/)\]

$\div$户 U2${\rm •}$ ， C (3) 绕平板的环量.

\noindent\textbf{7.19.} 图7.38 所示为拉瓦尔喷管流动.已知 A.f A T =2.005. 如= 1 X 105 Pa.

户。= 1. 36 X 106 Pa. (1)证明流动是管外膨胀; (2) 求出口压力户，及气流扩张角 $\alpha$. 图7.38 题 7.19 示意阁 飞dr气 /与市纣 图7.39 题 7.20 示意图

\noindent\textbf{7.20.} 图7.39 所示拉瓦尔喷管流动，已知 A ， /AT = 1 1. 91.$\rho$ol=7X10'Pa.

\[Pb=2. 26X 10"Pa. T 01 =500K.\]

(1)证明此流动管内有激波: (2) 求激波位置上的面积比 A ，/ A T ${\rm •}$ 及波前马赫数 Ma ，j'; (3) 求 Ma ，. T ,.

\noindent\textbf{7.21.} 超声速风洞内的流动如图 7.40 所示 .A 为风洞工作段面积.已知 :A=

1. 686An , A,=O. 822A ,A.=4. 130A.A=O. 2m2 ${\rm •}$ To =288K. 扣 =100kN/m 2 ${\rm •}$ 求 z (1)工作段马赫数 z (2) 扩压段喉部面积 A$\eta$ 的最小值; (3) 出口马赫数 p (4) 前室总压 PO; (5) 质量流量. A A. Ph 图 7.40 题 7.21 示意图

\section*{思考题}\addcontentsline{toc}{section}{思考题}

S7.1 绝热过程和等搞过程有什么区别? S7.2 流场中两侧垂(切〉向速度不同的平面称为法(切)向间断面，不可压缩流 动是否可能产生法向间断面?为什么?切向间断面两侧的压强是否相等?为什么? 不可压缩与可压缩流动，是否都可能产生切向间断面?为什么? S7.3 起点总馆和摘相等的两条流线，一条穿过激波，另一条不穿过激波，这两 条流线下游的总熔和摘是否相等?为什么? S7.4 气体穿过激波后，临界温度和临界压强是否改变?为什么? S7.5 理想完全气体分别以亚声速、超声速进入扩张通道时，是否可能产生激 波?为什么?如果进入收缩通道呢?
